\documentclass[11pt]{article}
\usepackage[utf8]{inputenc}
\usepackage[T1]{fontenc}
\usepackage{geometry}
\usepackage{amsmath}
\usepackage{amsfonts}
\usepackage{amssymb}
\usepackage{graphicx}
% \usepackage{natbib}
\usepackage{url}
\usepackage{hyperref}
\usepackage{titlesec}
\usepackage{enumitem}
\usepackage{wrapfig}
\usepackage[backend=biber]{biblatex}
\addbibresource{ref.bib}

\usepackage{booktabs}
\usepackage{array}
\usepackage{caption}
\usepackage{subcaption}
\usepackage{lmodern}
\usepackage{microtype}
\usepackage{threeparttable}   % for threeparttable and tablenotes
\captionsetup{font=scriptsize, labelfont=bf, position=top, skip=3pt}
\setlength{\tabcolsep}{4pt}
\renewcommand{\arraystretch}{1.05}

% Page layout
\geometry{
    left=0.8in,
    right=0.6in,
    top=0.6in,
    bottom=0.8in
}
\newcommand{\ts}{\mathrm{TS}}
\newcommand{\jsd}{\mathrm{JSD}}
\newcommand{\tsw}{\mathrm{TS_{wild}}}
\newcommand{\inm}{\mathrm{Infl}}
\newcommand{\adj}{\mathrm{Adj}}
\newcommand{\K}[1]{{\color{red}[[\textbf{K: }#1]]}}
\newcommand{\E}[1]{{\color{blue}#1}}
\newcommand{\facin}{f(\mathrm{act/inh})}
\newcommand{\ewnum}{\mathrm{EW_{num}}}
\newcommand{\tsten}{\mathrm{Toggle Switch (10N,100E)}}
\newcommand{\tstwo}{\mathrm{Toggle Switch (10N,80E)}}
\newcommand{\tsfour}{\mathrm{Toggle Switch (10N,60E)}}
\newcommand{\tssix}{\mathrm{Toggle Switch (10N,40E)}}
\newcommand{\tseight}{\mathrm{Toggle Switch (10N,20E)}}
\newcommand{\emttwosix}{\mathrm{EMT (23N,89E)}}
\newcommand{\emtrac}{\mathrm{EMT (17N,62E)}}
\newcommand{\gonad}{\mathrm{Gonadal (18N,74E)}}
\newcommand{\siltwo}{\mathrm{Silviera (18N,37E)}}
\newcommand{\sil}{\mathrm{Silviera (16N,30E)}}
\newcommand{\sclc}{\mathrm{SCLC (33N,357E)}}
\newcommand{\tennode}{\mathrm{NRF2 + OVOL (8N,19E)}}
\newcommand{\twenode}{\mathrm{NRF2 + GRHL2 + OVOL (11N,28E)}}
\newcommand{\beginsupplement}{%
        \setcounter{table}{0}
        \renewcommand{\thetable}{S\arabic{table}}%
        \setcounter{figure}{0}
        \renewcommand{\thefigure}{S\arabic{figure}}%
        \setcounter{section}{0}
        \renewcommand{\thesection}{S\arabic{section}}%
     }

%%%%%%%%%%%%%%%%%%%%%%%%%%%%%%%%%%%%%%%%%%%%%%%%%%
\title{Regulatory stochasticity drives opposing phenotypic outcomes in cell-fate decision networks\\

}

\author{Abhay Gupta, Lakshmi Malvadi, Prakash Kulkarni, Ravi Salgia, Herbert Levine, \\ Mohit Kumar Jolly, Kishore Hari}

\begin{document}

\maketitle

\begin{abstract}
Gene regulatory networks (GRNs) underlying cell-fate decisions are conventionally modelled as deterministic systems with fixed interaction parameters, yet the effective strength of any regulatory interaction fluctuates continuously in vivo. Intrinsically disordered regions in transcription factors, non-specific genomic binding sites, and stochastic enhancer dynamics collectively introduce temporal variability in the fold-change parameters that govern regulatory efficacy --- a regulatory noise that is distinct from, and largely independent of, classical copy-number noise. How such regulatory noise shapes the emergent phenotypic landscape of cell-fate circuits remains poorly understood. Here, we systematically evaluate the effects of three modes of regulatory noise --- additive (state-independent), multiplicative (state-dependent), and stationary (mean-reverting) --- on GRN motifs underlying cell-fate decisions. Using RACIPE-generated parameter ensembles, we simulate toggle switches, toggle switches with self-activation, and toggle triads under each noise mode across a wide range of noise magnitudes, and extend the analysis to larger cell-fate networks governing epithelial--mesenchymal transition and gonadal cell-fate determination, all using a Boolean framework. Additive noise consistently drives networks away from terminal attractors toward hybrid co-expression states, with the effect scaling with noise intensity and being strongest in bistable parameter regimes and sparser network architectures. Multiplicative noise, by contrast, reinforces inhibitory interactions over time and stabilises terminal phenotypes across all topologies studied. Stationary noise produces minimal disruption to the phenotypic landscape, consistent with a normal physiological baseline. These results establish that the type of regulatory noise, not only its magnitude, is a key determinant of phenotypic plasticity in cell-fate circuits, and suggest that additive-like noise --- plausibly arising from IDR conformational dynamics or open chromatin in stem cells --- may be a non-genetic driver of hybrid state enrichment and intra-tumour heterogeneity.
\end{abstract}

\section{Introduction}\label{sec:Intro}

GRNs are biological regulatory circuits that constitute a complex collection of molecular regulators including transcription factors (TFs), signaling molecules, and DNA segments that interact to regulate gene expression. They determine the spatiotemporal patterns of gene expression involved in determining cellular development, identity, and environmental responses. They also play a central role in determining the organism’s body plan, which is central to evolutionary developmental biology \cite{Davidson2005, Davidson2006, Maizels2026}. GRNs are conventionally modelled as deterministic systems with fixed interaction parameters, although it is clear that the effective strength of any regulatory interaction is subject to continuous fluctuation in vivo. For example, an overwhelming fraction of TFs are intrinsically disordered proteins (IDPs) with significant disordered regions . These disordered regions mediate conformationally heterogeneous, post-translational modification (PTM)-sensitive contacts with DNA and co-factors, causing regulatory affinities to vary dynamically rather than assume a fixed value \cite{Ferrie2022Structure-functionRegulation, Holehouse2023TheRegions}. Considering that 80-90\% of TFs are IDPs \cite{Staby2017}, understanding IDP conformational noise is paramount for understanding cell-fate specification. Although the individual conformers in an IDP ensemble are in fast exchange, certain conformations are stabilized by specific PTMs, especially phosphorylation, that exist with half-lives that range from minutes to hours contributing to ‘conformational noise’ \cite{Kulkarni2022, Kulkarni2019}. Non-specific TF binding to genomic decoy sites further modulates the effective regulatory signal in an affinity- and abundance-dependent manner \cite{Dey2020EnhancementSites, Burger2012InfluenceExpression}. Enhancer dynamics and transcriptional bursting add a third layer, coupling stochastic TF occupancy at cis-regulatory elements directly to fluctuations in target gene activation \cite{Raj2006StochasticCells, Fukaya2016EnhancerBursting}. The availability of TF binding sites can also vary over the course of transcription, due to noise in chromatin folding states \cite{Gamble2025}. Collectively, these mechanisms argue that the parameters governing regulatory interactions — rather than being static — are themselves stochastic quantities, and the consequences of such ''regulatory noise", i.e., temporal stochasticity in regulatory interaction parameters, for the dynamics of multistable GRN motifs underlying cell-fate decisions remain poorly understood.


This question is especially pressing in the context of cancer, where genetic instability, aberrant chromatin remodelling, and rewired signalling cascades collectively amplify both intrinsic gene expression noise and extrinsic cell-to-cell variability \cite{Brock2009Non-geneticTumours, Landau2014LocallyLeukemia, Hanahan2026HallmarksBeyond}. Non-genetic heterogeneity — stable phenotypic variation within a clonal population arising from gene expression noise and network multi-stability — is now recognised as a major contributor to drug tolerance and tumour progression \cite{Hanahan2026HallmarksBeyond, Franca2024CellularContinuum, Shaffer2017RareResistance}. In this setting, the parameters of regulatory interactions which govern the dynamics of non-genetic heterogeneity cannot be considered fixed: PTM landscapes, co-factor availability, and chromatin accessibility to TFs all fluctuate, effectively introducing stochasticity into the fold-change parameters that govern network behaviour. 


A natural starting point for studying how regulatory noise propagates through GRNs is the toggle switch — a simple mutual repression circuit between two master regulators that is among the most commonly occurring and extensively characterized motifs for binary cell-fate decisions \cite{gardner2000construction}. The toggle switch exhibits bistability under a broad range of kinetic parameters, and noise-induced switching between its two attractors has been studied in detail in both synthetic and natural contexts \cite{tian2006stochastic}. Many biological decisions, however, are not binary: T helper cell differentiation, haematopoietic lineage commitment, and epithelial–mesenchymal–hybrid state transitions all involve three or more stable phenotypes. Addition of self-activation to the toggle switch – TSSA (toggle switch with self-activation) stabilizes the intermediate phenotype – corresponding to hybrid phenotypes in EMT and progenitor states in differentiation cascades – where both nodes are expressed simultaneously \cite{Lu2013}. Three nodes connected to each other via a toggle switch, forming the toggle triad (TT) enables three distinct cell states, each defined by one node being highly expressed and the other two repressed \cite{doi:10.1098/rsif.2020.0631}. We will have occasion to study the 4-node toggle tetramer (see below). Importantly, although GRNs underlying cell-fate decision systems often involve a larger number of TFs than accounted for in these motifs, TFs corresponding to a phenotype often form teams such that nodes belonging to the same team activate each other. These teams mutually inhibit each other, thereby forming higher order forms of the motifs listed above. These higher order motifs have similar qualitative behavior as their small motif counterparts such as a two team network promoting bistability with a small fraction of intermediate states, similar to TSSA. 


The effect of stochasticity on toggle switch dynamics has been studied extensively, though almost exclusively through noise applied to molecular copy numbers rather than interaction parameters. Early work established that intrinsic noise — arising from the discreteness of protein and mRNA molecules — can drive spontaneous switching between the attractors of a multistable circuit, with switching rates that decrease exponentially with copy number and increase with noise magnitude \cite{tian2006stochastic, Wang2007Noise-inducedSwitch}. Analytical treatments using the Chemical Langevin Equation and mean first-passage time theory have characterized how the depth and asymmetry of the potential wells governing each attractor determine residence times and transition rates \cite{Wang2010TheDifferentiation, Strasser2012StabilityExpression, Xu2014SwitchSwitch}. A key finding across these studies is that noise type — not only noise magnitude — shapes attractor occupancy: additive and multiplicative noise applied to the same deterministic system produce different effective potential landscapes and therefore different stationary probability distributions \cite{Jaruszewicz2013ToggleGene}. Extending beyond Gaussian noise, Lévy noise models have shown that heavy-tailed fluctuations can induce coherent switching at substantially lower intensities than Gaussian noise requires, by enabling large jumps that entirely overcome the potential barrier \cite{Xu2016TheNoise}. At the network level, toggle switch dynamics have been studied under coupled and multicellular settings, where intracellular multiplicative noise can induce synchronised switching across cell populations via quorum-sensing pathways \cite{Wang2007Noise-inducedSwitch}. More recently, the sRACIPE framework demonstrated that expression-level (state) noise merges bistable clusters at high intensity, while static parametric variation across the model ensemble instead broadens clusters without destabilizing them \cite{Kohar2018RoleCircuits}. In both cases, however, the kinetic parameters themselves remain fixed once sampled; neither treats the parameter as a time-evolving stochastic process, which is the gap that our regulatory noise framework addresses. Stochastic analysis of the toggle triad and related tristable motifs has received comparatively little attention, with existing studies focused primarily on deterministic bifurcation structure and attractor geometry \cite{doi:10.1098/rsif.2020.0631, Duddu2025DistinctNetworks} rather than noise-driven state transitions. Whether the richer transition topology of the toggle triad — permitting direct switching between all three attractor pairs — confers qualitatively different susceptibility to regulatory noise compared to the toggle switch remains an open question.


In this manuscript, we study how regulatory noise – that is, the noise of parameters governing the interactions in GRNs – shapes the dynamics of these motifs systematically. We focus on the fold-change parameter \cite{Huang2017}, which describes the change in the production rate of a target node due to an upstream transcription factor. We couple stochastic fold-change perturbations with the Random Circuit Perturbation (RACIPE) framework \cite{Huang2017}. RACIPE generates large ensembles of ODE models with randomly sampled kinetic parameters, enabling topology-level inferences that are robust to parameter uncertainty. By superimposing continuous stochastic noise on the fold-change parameters of these ensembles, we successfully isolated the effect of regulatory noise from parameter-set-specific idiosyncrasies and characterized how mean residence times and switching frequencies depend on noise magnitude, noise mode, and attractor stability class. Our results reveal that regulatory noise is sufficient to induce state transitions that are inaccessible in the deterministic limit, and that bistable and tristable parameter regimes differ qualitatively in their susceptibility to such noise — with implications for understanding phenotypic plasticity in both normal development and disease.

\section{Results}

\subsection{Simulating regulatory noise in GRN motifs underlying cell-fate decision}


\begin{figure}
    \centering
    \includegraphics[width=\linewidth,height=0.7\textheight,keepaspectratio]{Figs/Fig1.jpg}
    \caption{
    \textbf{RACIPE network topologies and the stochastic simulation framework.}. 
    \textbf{(A)} Network topology diagrams for the four base circuits used throughout this study: (i) TS (toggle switch), (ii) TSSA (toggle switch with self-activation), (iii) TT (toggle triad), (iv) TTSA (toggle triad with self-activation). Red circles denote nodes; bar-headed edges denote inhibitory interactions; curved arrows denote self-activation.
    \textbf{(B)} (i) Schematic of the RACIPE parameter-sampling framework: production, degradation, and interaction-strength parameters are randomly sampled within a bounded parameter space, generating an ensemble of models with varying dynamical behavior. (ii) Steady-state expression of node A vs. node B (log scale) across the RACIPE parameter ensemble for the TS network. Dashed red lines mark each gene's discretization threshold (mean of its expression in log scale across all steady states); quadrants are labeled with the corresponding discrete state, (A, B).
    \textbf{(C)} Frequency of the 10 most frequent discrete attractors found by RACIPE for each network (TS, TSSA, TT, TTSA), expressed as the fraction of parameter sets that map onto these attractors. States are represented in short-hand form, with the key given inside the plot. Compound labels (e.g.\ ``A\_B'') denote multistable parameter sets. Because TT and TTSA support many more distinct multistable combinations than TS/TSSA, the top-10 bars shown capture only part of the full ensemble: 100\%, 99.2\%, 77.5\%, and 46.2\% of all parameter sets for TS, TSSA, TT, and TTSA respectively (out of 9, 12, 97, and 119 total distinct attractor combinations found).
    \textbf{(D)} Simulated trajectories of the fold-change parameter $\lambda(t)$ under the three stochastic noise-update rules used in this study (Stationary, Multiplicative, Additive), shown for an inhibitory-edge-range parameter ($\lambda \in [0,1]$, $\lambda_0 = 0.2$). Thin black dashed lines show 5 representative individual trajectories (sampled once per time unit for legibility); the colored line and shaded ribbon show the mean $\pm$ SD across 100 simulated trajectories.}
    \label{fig:1}
\end{figure}

We begin by analyzing the binary and tertiary cell-fate decision motifs --- the toggle switch (TS) and the toggle triad (TT) and their extensions that include self-activation edges (TSSA and TTSA). A toggle switch consists of two mutually inhibiting transcription factors (Figure~\ref{fig:1}A(i-ii)), whereas a toggle triad consists of three mutually inhibiting transcription factors (Figure~\ref{fig:1}A(iii-iv)). TS represents the core motif underlying binary cell-fate decision networks. The phenotypic landscape emergent from TS is dominated by two states with mutually exclusive expression patterns - either node has low expression while the other is high - $(0,1)$ or $(1,0)$. Under appropriate parametric conditions, the toggle switch also exhibits bistability, enabling transitions between $(1,0)$ and $(0,1)$ states \cite{gardner2000construction}. The addition of self-activation (TSSA) stabilizes the $(1,1)$ state in which both nodes are co-expressed. This co-expression state, often referred to as a "hybrid" state has been associated with stemness in development \cite{Roeder2006TowardsPU.1} and with high plasticity and pro-metastatic features in cancer \cite{Pastushenko2019EMTMetastasis}.

While binary cell-fate decision systems dominate biology owing to their stability and robustness to internal and external perturbations \cite{harlapur2022functional, Shyam2025MutuallyNetworks}, some instances of tertiary cell-fate systems --- such as T-cell differentiation --- employ a toggle triad architecture. A toggle triad consists of three mutually inhibiting transcription factors (Figure~\ref{fig:1}A(iii)) and permits three differentiated, single-high states where one of the three nodes has a high expression while the other two have low expression ($(0,0,1)$, $(1,0,0)$, and $(0,1,0)$) as well as three hybrid, double-high states with two out of three nodes having high expression (e.g.\ $(1,1,0)$). Parameter sweeps of the toggle triad result predominantly in monostable and bistable configurations, with a relatively smaller fraction exhibiting tristability. Single network architectures that allow the co-existence of four or more mutually exclusive states are not currently known.

We model these networks using a Hill-equation kinetics-based ODE system (Equation~\ref{eq:ode}, \ref{eq:hill}, Section~\ref{sec:RACIPE}). 
\begin{equation}
\frac{dX}{dt} = g_X \prod_{i \in P} \frac{H(Y_i, Y_{i0}, n_{Y_i}, \lambda_{Y_i})}{\lambda_{Y_i}} \prod_{i \in N} H(Y_i, Y_{i0}, n_{Y_i}, \lambda_{Y_i}) - d_X X
\label{eq:ode}
\end{equation}
\begin{equation}
H(Y_i, Y_{i0}, n_{Y_i}, \lambda_{Y_i}) = \lambda_{Y_i} \frac{Y_i^{\,n_{Y_i}}}{Y_{i0}^{\,n_{Y_i}} + Y_i^{\,n_{Y_i}}} + \frac{Y_{i0}^{\,n_{Y_i}}}{Y_{i0}^{\,n_{Y_i}} + Y_i^{\,n_{Y_i}}}
\label{eq:hill}
\end{equation}
Briefly, these equations combine transcription and translation into a single production term and a first-order degradation term representing mRNA and protein turnover. The production term consists of a rate constant ($g_X$) representing polymerization of RNA and protein, together with a combination of shifted Hill functions that formalize transcriptional regulation by input nodes. Input nodes change the production rate via the fold-change parameter $\lambda$: for inhibitory inputs, $\lambda < 1$, so that the presence of the regulator reduces the production rate; for activating inputs, $\lambda > 1$. The $\lambda$ of inhibitory links is bound between 0.01 (strongest inhibition) and 1 (no inhibition) and that of the activatory links is bound between 100.0 (strongest activation) and 1.0 (no activation). 
%The shifted Hill functions are formulated such that the maximum production rate of any node is always $g_X$, making the maximum achievable expression $g_X/d_X$.  The fold-change parameter thus represents the regulatory efficacy of an input node, making it a natural target for studying structural stochasticity. 

We first simulated these networks deterministically using RACIPE~\cite{Huang2017}, a tool that allows us to explore the different configurations that the emergent phenotypic landscape can take across a large parameter space. RACIPE generates a large ensemble of ODE models by sampling multiple parameter sets (Figure~\ref{fig:1}B) and simulates each ODE model from multiple initial conditions until steady states are reached, resulting in an ensemble of steady states. We discretized expression of each node in each steady state into high (1) or low (0) based on the node's log-average expression across the ensemble of steady states (Figure~\ref{fig:1S1}B).  Threshold values for each network are given in Table~\ref{tb:Thresholds}. The models generated by RACIPE can be classified as monostable or bistable based on the number of discrete steady states they support, and further subdivided by the specific attractor configuration (see Table~\ref{tb:Attractors}). For the toggle switch, the dominant attractors were $(0,1)$, $(1,0)$, and the bistable $(0,1)\text{--}(1,0)$, as expected (Figure~\ref{fig:1}C, i). Although a small fraction of models do permit $(1,1)$ as a steady state, the corresponding fold-change parameters are always close to 1, rendering inhibitions ineffective~\cite{hari_assessing_2023}. TSSA enhances bistability in the parameter space, with the other dominant attractors being $(0,1)$ and $(1,0)$. Additionally, TSSA also supports tristability with $(0,1)$ and $(1,0)$ and either $(0, 0)$ or $(1, 1)$. The $(0, 0)$ state was accessed by a very small number of initial conditions, while $(1, 1)$ was more accessible (Figure~\ref{fig:1}C, ii). TT and TTSA support single-high (one of the nodes has a high expression - $(0, 0, 1)$ etc.) and double-high (two of the nodes have high expression - $(0, 1, 1)$ etc.) nearly equally. TTSA enhances multistability over TT, promoting a tristable attractor with three single-high states (Figure~\ref{fig:1}C, iii, iv).

We explored three noise modes (Section~\ref{sec:noise}) described by the equations below: 
\begin{equation}
\begin{aligned}
\text{Additive:} \quad & \lambda(t+dt) = \lambda(t) + \mathcal{N}(0, \sigma_{\text{eff}}^2) \\
\text{Multiplicative:} \quad & \lambda(t+dt) = \lambda(t)\left(1 + \mathcal{N}(0, \sigma_{\text{eff}}^2)\right) \\
\text{Stationary:} \quad & \lambda(t+dt) = \lambda(0) + \mathcal{N}(0, \sigma_{\text{eff}}^2)
\label{eq:lambda}
\end{aligned}
\end{equation}

Networks were simulated at eight levels of $\sigma$ 0.0, 0.001, 0.005, 0.01, 0.05, 0.1, 0.5 and 1.0 scaled by the upper bound of $\lambda$ (1.0 for inhibitions and 100.0 for activations). The choices of $\sigma$ were made to ensure that low to moderate noise region has ample representation, as that region likely represents the underlying biology better than unrealistically high noises that span 50 or 100\% of the parameter range. The higher $\sigma$ values serve a more theoretical purpose in helping us understand the effect of stochastic simulations. The three noise modes produce qualitatively distinct trajectories of $\lambda$ over time . Stationary noise models normally distributed fluctuations around the initial value of $\lambda$ (Figure~\ref{fig:1}C, i, \ref{fig:1S1}A, i). Additive noise models a constrained random walk of $\lambda$ bounded by the corresponding limits for activating and inhibiting $\lambda$ values. The resultant stationary distribution has a mean of 0.5 for inhibitory links and 50 for activatory links (Figure~\ref{fig:1}C, ii, \ref{fig:1S1}A, ii). Multiplicative noise scales $\sigma$ by the current $\lambda$ value, such that lower noise values show lower deviation and higher noise values show higher deviation (Figure~\ref{fig:1}C, iii, \ref{fig:1S1}A, iii). The consequences of these distinct $\lambda$ trajectories for network state dynamics are explored in the following sections.


We sampled monostable and bistable models separately from RACIPE, weighting by the relative frequencies of each attractor configuration. We then performed stochastic simulations of these networks by periodically perturbing $\lambda$. For each parameter set, we generated 100 stochastic trajectories starting from randomly chosen initial conditions. We use the RACIPE ensemble-derived thresholds to discretize the stochastic trajectories as well (Figure~\ref{fig:1}D). Based on the number of discretized steady states allowed by a given parameter set in the deterministic simulations, each parameter set was characterized as monostable (single steady state), bistable, tristable and so on. For each network, we sampled 100 parameter sets per stability class and for each parameter set, we performed 100 stochastic simulations each to estimate the statistics. 



\subsection{High levels of stationary noise pushes the system towards hybrid states}


\begin{figure}
    \centering
    \includegraphics[width=\linewidth]{Figs/Fig2.jpg}
    \caption{\textbf{Effect of stationary noise on toggle switch and toggle triad networks}. 
    \textbf{(A)} Representative single-parameter trajectory for the TS network: (i) node A and (ii) node B expression over time (log$_2$), with the deterministic (noise-free) segment at negative time and the stochastic segment at positive time. Grey lines show all sampled initial conditions; the black line highlights one representative trajectory (the one whose switch count is closest to the population median); dashed horizontal line marks each node's discretization threshold $\theta$.
    \textbf{(B)} Discrete-state trace of the same parameter, colored by state for the highlighted trajectory (grey for all other initial conditions). Crosshatched region marks the discarded burn-in window of the stochastic segment. Corner text reports, per state, the fraction of time the highlighted trajectory (``Traj.\ residence time'') and the full population of initial conditions (``MRT'', mean residence time) spend in each discrete state.
    \textbf{(C)} MRT vs.\ noise level for the TS network (all parameter sets), faceted by discrete state; individual parameter-set trajectories (grey) with population mean (dashed black).
    \textbf{(D--F)} MRT vs.\ noise level for TSSA (D), TT (E), and TTSA (F) bistable parameter sets, colored by state (top 4 states shown for TSSA, top 7 for TT/TTSA), mean $\pm$ SD across parameter sets. Caption below each panel reports, per state, whether its MRT trajectory changes significantly across noise levels (Friedman test, blocked by parameter set; ``*'' $p-value<0.05$, ``ns'' not significant).
}
    \label{fig:2}
\end{figure}

Stationary noise mimics the classical assumption underlying GRN models, that the fitted parameters are the mean of a distribution that does not significantly change the emergent phenotypic landscape from what the mean alone predicts. Thus, we started with analyzing the stochastic simulations under stationary noise to evaluate the validity of deterministic representation of regulatory parameters. We analyzed the stochastic trajectories by first discretizing them using the RACIPE ensemble threshold values (Figure~\ref{fig:2}A). The resultant discrete state trajectories are then used to calculate the residence time of each state, defined as the fraction of time steps a given trajectory sees that state. An average of the residence time over 100 trajectories, the mean residence time (MRT) is used to evaluate the occurrence of a state for a given parameter set under a given noise type and noise level (Figure~\ref{fig:2}B). Note that we use the same MRT calculation for deterministic simulations here on (e.g., $sigma = 0$). As deterministic simulations reach steady state by the burnin period (shaded in criss-crossed lines in Figure~\ref{fig:2}B), the MRT measures the relative size of the basin of attraction. For monostable parameters, MRT in deterministic simulations will always be 1, while multistable parameters will see fractional MRT corresponding to the basin size of the specific steady state in the attractor. 

As expected, the dominant states for a toggle switch remain as $(0, 1)$ and $(1, 0)$ for low to moderate $\sigma$ values under stationary noise. However, as the $\sigma$ increases to 0.5 and above (i.e., the noise scales as 50\% of the allowed range of $\lambda$), we see a drop in the MRT of the $(0,1)$ and (1, 0) states and a corresponding increase in the MRT of the (1, 1) state - which represents a co-expression of both nodes that inhibit each other (Figure~\ref{fig:2}C). Note that a similar increase is not seen in the MRT of the (0, 0) state (Figure~\ref{fig:2S1}A), an asymmetric outcome despite the symmetry of the stationary noise. We further studied TSSA, TT and TTSA under stationary noise and found them behaving similarly. TSSA, in absence of noise enables bistability in a larger fraction of parameter space than TS ($\approx 20\%$ in TS vs $>50\%$ in TSSA). TSSA also enables ~14\% of the state space to be tristable, with access to either $(0, 0)$ or the $(1, 1)$ state in addition to the dominant $(0, 1)$ and $(1, 0)$ states. However, the basin size of $(0, 0)$ state, is negligibly small, making $(0, 1) - (1, 0)-(1, 1)$ the only real tristable configuration (Figure~\ref{fig:2S1}B - tristable parameters at $\sigma = 0$). Similar to TS, TSSA maintains the characteristics of its phenotypic landscape under stationary noise - dominant $(0,1) and (1,0)$ states in monostable and bistable parameter sets, and dominant $(1, 1)$ state in tristable parameter sets for realistic values of noise, i.e., $\sigma \leq 0.1$ (Figure~\ref{fig:2}D). For higher $\sigma$, we see a dominance of $(1, 1)$ state. 

Toggle triad (TT) and toggle triad with self-activation (TTSA) dominantly enable two classes of steady states - single-high states ($(0, 0, 1); (0, 1, 0)$ and $(1, 0, 0)$ and double-high states ($(0, 1, 1); (1, 0, 1)$ and $(1, 1, 0)$). Previous studies have shown that toggle triad motif is more susceptible to various perturbations due to the inherent frustration in the network~\cite{harlapur2022functional, doi:10.1098/rsif.2020.0631}. Here, we see that the MRT distribution of TT starts to change for $\sigma > 0.01$, showing a higher sensitivity to noise than TS and TSSA motifs (Figure~\ref{fig:2}E, F, Figure~\ref{fig:2S1}C-D). At $0.01 \leq \sigma \leq 0.1$, the MRT of single-high states starts to go down and that of double-high states starts to come up, still qualitatively retaining the properties of the phenotypic landscape. However, $\sigma \geq 0.5$ leads to an increase in the MRT of the all-high state ($(1, 1, 1)$) in both TT and TTSA, similar to that of TS and TSSA, with no emergence of the all-low state ($(0, 0, 0)$). Overall, we find that the phenotypic landscape of both binary and tertiary cell-fate GRNs remains robust to stationary noise of low to moderate levels of $\sigma$. The response to larger noise can be better understood in our analysis of additive noise in the next section. 


\subsection{Additive noise reshapes the phenotypic landscape toward high-expression states, modulated by reachability and threshold margin}

Unlike stationary noise, additive noise starts to change the phenotypic landscape even at low $\sigma$ values. For toggle switch across a majority of the parameters, we see a decrease in the MRT of single-high states ($(0, 1)$ and $(1, 0)$) and an increase in the hybrid all-high ($(1, 1)$) state MRT (Figure~\ref{fig:3}A). The mean MRT of all-high state crosses 0.5 for $\sigma > 0.001$ (0.1\% of the lambda range) and stabilizes at 0.75 for $\sigma > 0.1$. Given the biological relevance of hybrid states showing co-expression of all nodes in developmental and pathological contexts, this shift in the phenotypic landscape of cell-fate decision motifs is a very significant result. 

\begin{figure}
    \centering
    \includegraphics[width=\linewidth]{Figs/Fig3.jpg}
    \caption{\textbf{Effect of Additive noise on the phenotypic landscape of cell-fate decision networks}
    \textbf{(A)} MRT of $(1, 1)$ and $(1, 0)$ states vs.\ noise level for the TS network, individual parameter-set trajectories (grey) with population mean (dashed).
    \textbf{(B)} MRT boxplots for the TS network, monostable vs.\ bistable parameter sets, all-high vs.\ single-high states, with paired Wilcoxon significance brackets; ``*'' $p-value<0.05$, ``ns'' not significant. 
    \textbf{(C)} Reachability scatter for TS monostable parameter sets at $\sigma = 0.1$: each node's (max expression $-$ threshold) margin on the x/y axes, colored by MRT of the (1,1) state.
    \textbf{(D)} Spearman correlation between the minimum reachability margin across nodes and MRT of the (1,1) state, across four noise levels (0.001, 0.01, 0.1, 1), TS monostable parameter sets.
    \textbf{(E)} MRT vs.\ noise level for TT and TTSA, single-high $\to$ double-high $\to$ all-high state classes shown as Mean $\pm$ SD across all parameters, with paired Wilcoxon significance brackets between adjacent classes; ``*'' $p-value<0.05$, ``ns'' not significant. 
    \textbf{(F)} Heatmap of the Spearman correlation between minimum reachability margin and MRT of the all-high state, across networks (TS, TSSA, TT, TTSA), noise levels, and stability class (monostable vs.\ multistable); ``*'' $p-value<0.05$, ``ns'' not significant.
}
    \label{fig:3}
\end{figure}

Interestingly, we note that while most parameters gravitate towards a high MRT for the all-high state for $\sigma \geq 0.01$, a smaller fraction of the parameters also show a significantly lower MRT of all-high state (Figure~\ref{fig:3}A). When we separately studied the parameter sets based on their multistability characteristics under deterministic conditions, we found that monostable states always had a lower MRT for all-high states under additive noise as compared to that of multistable parameter sets (Figure~\ref{fig:3}B). Given how readily bistable parameters converge to the all-high state, we hypothesized that the reduced MRT for monostable states could be because the all-high state is not accessible to some monostable parameter sets.  The upper limit of node expression in the ODE models we simulated is the ratio of the corresponding production and degradation rates. As we use a threshold derived from the ensemble steady state expression levels, it is possible that some parameter sets are sampled such that the upper limit of expression for one of the nodes defined by the production and degradation rates is lesser than the threshold. Since the bistable parameters constitute of two single-high states, bistable  parameter sets do not have this shortcoming. Thus, we wanted to understand whether it really is the case that monostable parameters that cannot access the all-high state are sampled, and in turn, if the access to all-high state is sufficient for the additive noise to push the system towards all-high state. 

Thus, we plotted the sampled parameter sets on the axes of "reachability" defined by the difference between the upper limit of a node's expression level and the threshold used the discretize that node. Reachability was satisfied across all sampled parameters — the production/degradation-derived expression ceiling exceeded threshold in every case. The observed variation in all-high MRT among monostable parameter sets therefore cannot be attributed to reachability itself (Figure~\ref{fig:3}C, \ref{fig:3S1}A). However, we see that the MRT of the all-high state increases with increasing distance between the maximum allowed value of the nodes and the corresponding thresholds. Across $\sigma$ values, the MRT of the all-high state correlated strongly positively with the smallest of the two difference values (Figure~\ref{fig:3}D). 

TSSA showed similar increase in hybrid (all-high) state MRT under additive noise. At low to moderate $\sigma$, TS and TSSA showed very similar MRT distributions. At high $\sigma$ values, TSSA monostable parameter sets showed a higher MRT of all-high states compared to TS (Figure~\ref{fig:3S1}B, i). This increase can be attributed to the inherent affinity TSSA has towards all-high state as discussed earlier. For TT and TTSA, we observed a transition of states from single-high to double-high to all-high, suggesting that while additive noise does push the system towards the all-high state, the transformation can be gradual as the network size increases (Figure~\ref{fig:3}E, \ref{fig:3S1}B, ii, iii). TT and TTSA have similar MRT distributions for low to moderate values of $\sigma$, with similar average MRT of all-high and double-high states. But at high values of $\sigma$, TTSA showed a higher MRT of the all-high state, with the double-high state average MRT approaching 0. TT on the other hand maintains a non-zero average MRT of double-high states even at higher $\sigma$ values (average MRT across all parameter sets - all-high 0.7 vs double-high 0.3 for TT vs all-high 0.9 and double-high 0.15 for TTSA). For all networks , we found a strong positive correlation between the difference quantity (max node expression minus threshold expression of the node) and the MRT of all-high states in monostable parameter sets. Mutlistable parameter sets showed relatively weaker correlations (Figure~\ref{fig:3}F). Thus, our analysis shows that additive noise increases the MRT of hybrid, all-high states in cell-fate decision networks. 

\subsection{Increased MRT of all-high states is caused by the $\lambda$ distribution and stochasticity}

Following the strong correlation we observed between parametrically determined upper bounds of node expression and the MRT of all-high state, we hypothesized that the drastic increase in the MRT of all-high states in comparison to deterministic simulations might also be explained by variation in the parameter space. Additive noise models a random walk in $\lambda$, bounded between 0.01 and 1 for inhibitory links and 1 and 100 for activating links. Any step crossing the boundary is absorbed at the boundary. For example, if $\lambda(t) + \eta(t) < 0.01$ then $\lambda(t+dt) = 0.01$. Due to the nature of random walk, the variance of lambda increases over time, unlike the stationary noise case. Furthermore, the absorbing boundaries increase the chances of $\lambda$ being at either of the borders than in the interior. The probability density of $\lambda$ at the borders increases with increasing $\sigma$, as at larger sigma values, there is a higher chance of taking larger $\lambda$ steps and correspondingly higher chances of hitting the boundaries (Figure ~\ref{fig:3b}A). Same boundary conditions are applied to other noise forms as well, and we see such boundary behavior in stationary noise at high $\sigma$, where $\lambda$ localizes at either of boundaries. While $\lambda$ closer to zero leads to a very strong inhibition by drastically reducing the net production rate, $\lambda$ closer to 1 leads to very weak inhibition and lets the production rate reach the highest values. Despite this symmetry in the emergent $\lambda$ distribution, the shift in the landscape is asymmetric towards all-high states. One potential reason could be the way $\lambda$ parameter is sampled in RACIPE for inhibitory interactions. The fold-change parameter is first sampled between 1 and 100 and then the inverse of that number is used for inhibitory links. Thus the distribution of $\lambda$ for inhibitory links is a log-uniform distribution tailed heavily towards lower values of $\lambda$ with a mean of $\approx0.05$. As the stationary distribution of $\lambda$ in additive noise has a mean of 0.5, inhibitory links experience a significant increase in the average $\lambda$, which in turn leads to weakening of the inhibition. Thus we wanted to know if this shift in $\lambda$ distribution is enough to explain the increase in the MRT of all-high states. For each parameter set used for stochastic simulations, we sampled $\lambda$ from the additive noise stationary distribution shown in (Figure ~\ref{fig:3b}A). For each such $\lambda$ combination, we simulated the networks deterministically ($\sigma = 0$) at multiple randomly sampled initial conditions and calculate the mean frequency of each discrete state across initial conditions, which is further averaged over different lambda values. While the original RACIPE parameters for TS and TT have a very small fraction of parameter sets with non-zero frequency of all-high states, majority of resampled simulations led to a non-zero frequency of all-high states, reaching up to 0.9 in TS and up to 0.75 in TT (Figure~\ref{fig:3b}B, \ref{fig:3bS1}A). TSSA and TTSA, which inherently stabilize the all-high states due to the self-activation links, resampled simulations still led to an increase in the frequency of all-high states for a majority of the parameter sets (Figure~\ref{fig:3b}B, \ref{fig:3bS1}A). Thus, the $\lambda$ distribution does contribute to increasing the MRT of all-high state. Interestingly, however, when we compared the MRT of all-high states in stochastic simulations and the resampled deterministic simulations, we found that the stochastic MRT was higher across a large fraction of parameters for all four networks (Figure~\ref{fig:3b}C, i, , \ref{fig:3bS1}B). Furthermore, a sigmoidal relationship emerges between the stochastic MRT and the deterministic resampled frequency of all-high state (Figure~\ref{fig:3b}C, ii, \ref{fig:3bS1}C). Below a threshold value of all-high state frequency under deterministic simulations with resampled $\lambda$, stochastic MRT is lower than or equal to deterministic frequency, and vice-versa (Figure ~\ref{fig:3b}C). The specific threshold MRT is network and $\sigma$ dependent (Figure~\ref{fig:3b}C, ii, \ref{fig:3bS1}C). Across all four networks and $\sigma$ values, the MRT of all-high states increased significantly from original to resampled deterministic to stochastic simulations (Figure ~\ref{fig:3b}D). We further tested the effect of stochasticity by reducing the frequency of update of $\lambda$, i.e., the $dt$ value in Equation~\ref{eq:lambda}, from 0.01 to 1.0 and 10.0 (Figure~\ref{fig:3bS3}). As the frequency of update decreased, the system had more time to relax towards a steady state, thereby losing the effects of stochasticity. Correspondingly, we see a reduction in the MRT of all-high state and increase in the MRT of the native states of these networks at lower $\sigma$ values. At higher sigma values, the noise compensates for the relaxation, pushing the systems back to hybrid states. These results indicate that, while the distribution of $\lambda$ does shift the landscape towards the all-high hybrid state, the stochastic simulations serve to further enhance this shift. 

\begin{figure}
    \centering
    \includegraphics[width=\linewidth]{Figs/Fig4.jpg}
    \caption{\textbf{The additive noise effect combines a $\lambda$-distribution shift with an additional stochastic contribution.}
    \textbf{(A)} Distribution of $\lambda$ (Hill-function threshold) values across the parameter ensemble, for (i) inhibitory and (ii) activation-type edges, at four noise levels (Additive noise).
    \textbf{(B)} Distribution of the resampling effect - difference in deterministic MRT after $\lambda$-resampling and the original point-$\lambda$ deterministic MRT for the all-high state, shown as violin plots across all noise levels, for TS, TSSA, TT, and TTSA.
    \textbf{(C)} (i) Same as B, but for Stochastic effect (stochastic MRT minus resampled-$\lambda$ deterministic MRT). (ii) Scatter plot of resampled-$\lambda$ deterministic vs.\ stochastic MRT of the all-high state at $\sigma = 0.1$, for TS and TT, illustrating the underlying sigmoidal relationship the violin's signed difference alone does not show. Each point is a parameter set.
    \textbf{(D)} All-high-state MRT (mean $\pm$ SD over all monostable and multistable parameter sets) for TS, TSSA, TT, and TTSA, comparing the original point-$\lambda$ deterministic run, the $\lambda$-resampled deterministic run, and the full stochastic simulation; faceted by stability class of the parameter sets (rows) and noise level (columns: 0.001, 0.01, 0.1), with paired Wilcoxon significance brackets; ``*'' $p-value<0.05$, ``ns'' not significant.}
    \label{fig:3b}
\end{figure}


To isolate the mechanism by which regulatory ($\lambda$) noise redistributes occupancy toward high-expression states, we analyzed the toggle switch in the Boolean limit of node expression and $\lambda$ distribution, where the exact stationary distribution can be computed in closed form for both the deterministic (fixed-$\lambda$) and stochastic (fluctuating-$\lambda$) cases (see Supplementary section \ref{sec:boolean}). We noticed a crucial asymmetry in the stochastic simulations that is absent in deterministic simulations. For an inhibitory edge, the shifted Hill function entering the production term is
\begin{equation}
  H^{S}(x,\lambda) \;=\; H^{-}(x) \;+\; \lambda\,H^{+}(x),
  \qquad H^{+}(x) \equiv 1 - H^{-}(x),
  \label{eq:hshift}
\end{equation}

Under the extreme case of $\sigma=1$, where the $\lambda$ walk is confined to the two boundaries of its range, $\lambda\in\{\lambda_{\min},\lambda_{\max}\}$ with $\lambda_{\min}\approx 0.01$ and $\lambda_{\max}=1$. Evaluating Eq.~\eqref{eq:hshift} at these two limits exposes a structural asymmetry:
\begin{align}
  \lambda = 1: \quad & H^{S}(x,1) = H^{-}(x)+H^{+}(x) = 1
    && \text{(independent of } x\text{)},\\
  \lambda \approx 0.01: \quad & H^{S}(x,0.01) \approx H^{-}(x)
    && \text{(depends on } x\text{)}.
\end{align}
The $\lambda=1$ boundary abolishes regulation \emph{unconditionally}: the target node is produced at its full rate regardless of the regulator's expression. The $\lambda\approx 0$ boundary imposes repression only \emph{conditionally}: it lowers the target's production rate only when the regulator $x$ is high, and has essentially no effect when the regulator is low ($H^{-}(x)\!\approx\!1$). An excursion toward de-repression ($\lambda=1$) therefore always ``fires,'' whereas an excursion toward repression fires only when the partner node happens to be expressed. This unconditional-up / conditional-down asymmetry is the microscopic origin of the enhancement. We see this clearly in the boolean approximation of node expression (i.e., each node can only take two expression levels - either 1 or 0). In the extreme $\sigma$ condition, resampled deterministic simulations can only converge to the all-high state if they start from the all-high state (a frequency of 0.25). However, in stochastic simulations, asymmetry in the response of shifted hill function ensures that a strongly repressed state is maintained only when the source node is on, or when both are off. This asymmetry leads to a higher probability (equivalent to MRT) of all-high state in the stationary distribution of the stochastic process. Even at lower values of $\sigma$, this asymmetry holds true by making the de-repressed state more easily accessible to the nodes of the network than the repressed states. Therefore, we conclude that the drift in the phenotypic landscape caused towards the all-high state caused by additive noise is a combined outcome of the shift in the parameter space and the asymmetry exploited by stochastic simulations. 


\subsection{Multiplicative noise drives the phenotypic landscape toward single-high states - an opposite effect as that of additive noise}


\begin{figure}
    \centering
    \includegraphics[width=\linewidth]{Figs/Fig5.jpg}
    \caption{Effect of multiplicative noise on the phenotypic landscape of cell-fate decision networks.
    \textbf{(A)} MRT vs.\ noise level for the TTSA network, single-high and double-high states.
    \textbf{(B)} Average MRT of single-high vs.\ double-high states vs.\ noise level, for TS, TSSA, TT, and TTSA.
    \textbf{(C)} Same, for the 4-node networks TT4 and TS4.
    \textbf{(D)} Stochastic effect (as in Fig.\ 4C(i)) for the single-high state, across noise levels, for TS, TSSA, TT, TTSA, TT4, and TT4SA.
}
    \label{fig:4}
\end{figure}

We then analyzed the effect of multiplicative noise in $\lambda$ on the phenotypic landscape. Briefly, multiplicative noise models a random walk in $\lambda$ with a state-dependent step size. For inhibitory interactions, the noise amplitude is smaller near $\lambda \approx 0$ and larger near $\lambda \approx 1$. This asymmetry makes $\lambda$ approach over $0$ over time, unlike additive noise where $\lambda$ visits both the boundaries and the intermediate values at intermediate $\sigma$ (Figure~\ref{fig:1}D). While this concentration of $\lambda$ near $0$ may seem like a deterministic behavior, stochasticity has a strong impact on the outcome of multiplicative noise, even more so than additive noise in certain cases as we will see below. 

Intuitively, multiplicative noise enhances repression in the edges of toggle switch, thereby driving the system strongly towards single-high states (Figure~\ref{fig:4S1}A, i). This does not affect the phenotypic landscape of TS, as single-high states were dominant to begin with (Figure~\ref{fig:1}C). TSSA also showed a similar gravitation towards single-high states despite TSSA also stabilizing the all-high state in absence of noise (Figure~\ref{fig:4S1}A, ii). The enhancement of single-high state by multiplicative noise is prominently seen in case of TT and TTSA. Both these networks have a significantly high frequency of double-high state along with single-high state in absence of noise (Figure~\ref{fig:1}C). In presence of multiplicative noise, both these networks converge to single-high states, achieving a high MRT (~1) by $\sigma = 0.01$ (Figure~\ref{fig:4S1}A, iii,~\ref{fig:4}A). Thus, consistently across the four networks studied so far, we found that multiplicative noise promotes the frequency of single-high states (Figure~\ref{fig:4}B). At very high $\sigma$ values, TTSA shows a decrease in the frequency of single-high states with a corresponding increase in double-high state MRT. At high $\sigma$, $\lambda$ starts to approach higher values purely based on the fluctuations (Figure~\ref{fig:4S1}B). However, the probability distribution of $\lambda$ still remains concentrated dominantly near the lower limit, resulting in an increase of double-high states only and not all-high states. To further validate this result, we simulated three new networks that are seen in cell-fate decision systems. The first network is an extension of TT network in the context of T cell differentiation, including the regulation of T-reg cells through the transcription factor Foxp3. The resultant four node network is referred to as toggle tetrahedron (TT4), and is composed of four nodes mutually inhibiting each other. Previous works have explored this network and its self-activation counterpart (TT4SA), our second 4 node network, as candidates for four-fate differentiation~\cite{Duddu2024MultistabilityDifferentiation}  (Figure~\ref{fig:4S1}C, i, ii). However, this and other candidate networks for multi-fate systems have so far failed to support the emergence of single-high states representative of fully differentiated cell-fates. TT4 and TT4SA, specifically support only double-high steady states (Figure~\ref{fig:4S1}D). Given the properties of multiplicative noise seen so far, we wanted to study if the double-high states of TT4 and TT4SA could also be driven towards single-high states. The third network is a combination of two mutually inhibiting toggle switches that we refer to as TS4, a representation of coupled cell-fate decision systems often seen in biology~\cite{pasani2020hybrid, hari2022emergent}. TS4 also supports double-high states (Figure~\ref{fig:4S1}D). However, unlike TT4 where all six possible double-high states have equal frequency, TS4 only supports a limited number of double-high states that represent the coexistence of the two binary cell-fate systems. Unlike sequential differentiation, coupled systems are expected to have co-existing phenotypes governed by the network structure. Thus, we asked whether multiplicative noise pushes these coexisting double-high states into single-high states, irrespective of the topology. Under additive noise, all three networks lead to a consistent increase in all-high state, despite the differences in the network structure (Figure~\ref{fig:4S2}A). Under multiplicative noise, TT4 saw an increase in the single-high state MRT with $\sigma$, followed by a decrease brought on due to the boundary sampling of $\lambda$ (Figure~\ref{fig:4}C, \ref{fig:4S2}B). However, TT4SA maintained a double-high dominant landscape (Figure~\ref{fig:4S2}B). TT4SA did show a decrease in double-high state MRT as multiplicative noise increased, with a corresponding increase in the MRT of single-high states. However, this change never put single-high states in a dominant position. Furthermore, TS4 also maintained a double-high dominant landscape, with no change in the MRT with $\sigma$, preserving the behavior of internal toggle switches. Thus, unlike additive noise, the effect of multiplicative noise on the phenotypic landscape is dependent on the network topology. 

Given the potentially deterministic nature of multiplicative noise, we tried to dissociate the effects of the $\lambda$ distribution from that of stochasticity in these simulations. Thus, we generated resampled deterministic simulations similar to that of Figure~\ref{fig:3b}C. For TS and TSSA, the difference between stochastic and resampled deterministic MRT was very low, as expected (Figure~\ref{fig:4}D). However, TT, TTSA and TT4 showed a high difference even at lower values of $\sigma$ (< $0.05$, Figure~\ref{fig:4}D). In comparison, additive noise simulations required higher  $\sigma$ ($\geq 0.1$) to achieve this level of difference (Figure~\ref{fig:3b}C). TT4SA and TS4 also showed minimal difference, in accordance with their MRT trends. These differences were mirrored by corresponding decreases in the MRTs of single-high states (Figure~\ref{fig:4S2}D, E). Interestingly, however, the MRTs of original deterministic simulations and resampled simulations were nearly identical (Figure~\ref{fig:4S2}F, G), indicating that this "fixation" towards terminal differentiation seen in multiplicative noise is entirely due to the stochasticity, unlike the additive noise where the $\lambda$ distributions also contributed to the increase in MRT of all-high states. We further tested this behavior by simulating these networks at larger dt values (Figure~\ref{fig:5S3}). Unlike additive noise, increasing dtdt dt did not reduce the excess of stochastic MRT over the resampled-deterministic MRT for multiplicative noise. Since larger dtdt dt gives node expression more time to relax toward the steady state set by the current $\lambda$ before the next update, this indicates that multiplicative noise's dynamical contribution — unlike additive noise's — is not primarily a consequence of expression lagging behind a fast-moving $\lambda$. The specific mechanism by which individual stochastic realizations of $\lambda$ produce this persistent excess remains to be characterized, but our results indicate it is distinct from the timescale-separation (non-adiabatic) effect identified for additive noise.

% We wanted to further understand the nature of these transitions - which single-high states do the hybrid double-high states converge to? Can multiple single-high states be achieved from a given double-high state - i.e., if we think of the double-high state as a progenitor, can it always differentiate into multiple terminal phenotypes? We chose to do this analysis on TT4, as this network has the highest frequency of double-high and lowest frequency of single-high states, thereby giving us the cleanest signal with a large dataset. TT4 predominantly allows monostable and bistable double-high steady states. We first analyzed the monostable steady states, and found that the single-high states emergent from the application of multiplicative noise were always the closest single-high states from the double-high states, achieved by pushing the expression of one of the high expression nodes to low. Bistable states of TT4 are composed of "mirror" states, i.e., different pairs of nodes have high expression in both the states. Therefore, the bistable double-high states have access to all four of the single-high states. We further find that while all these parmeters converge to single-high states upon application of appropriate multiplicative noise, a large fraction of parameters a small fraction of parameter sets (~19\% for TT, 33\% for TT4) converge to two single-high states in different instances of simulations. We further studied the transitions observed in multiplicative noise simulations and found that the only transitions are between the double-high states and single-high states, and never between the two single-high states. Furthermore the number of transitions observed are also very small (1 in every x time steps). Since we know that the initial time steps are the most noisy in multiplicative noise, most of these transitions would have happened initially, resulting in stable terminal states achieved after long-term application of multiplicative noise. We were able to confirm these results for TT, TTSA and TT4SA. Therefore, we concluded that multiplicative noise can push the landscape of cell-fate decision systems towards single-high, terminally differentiated states. 


\subsection{Effects of regulatory noise on the phenotypic landscape scale to larger biological networks}

While the toggle switch and TSSA are well-studied representatives of binary cell-fate decision systems, the regulatory circuits underlying real cell-fate transitions are rarely as small or as isolated. Phenotypes emerge from coordinated expression across tens to hundreds of genes, many of which function as transcription factors that individually regulate transitions between phenotypes. It is therefore important to ask whether the noise-driven dynamics observed in small motifs generalise to larger, more complex GRNs. To address this, we asked whether biological GRNs governing two well-studied cell-fate decisions -- Epithelial mesenchymal plasticity~\cite{hari2022landscape} and Gonadal cell-fate decision~\cite{rios2015boolean} -- behave consistently.

% \subsubsection{Artificial teamed networks}

% Large-scale GRNs underlying binary cell-fate decisions share a common coarse-grained structure: nodes can be partitioned into two ``teams'', with intra-team interactions predominantly activating and inter-team interactions predominantly inhibitory (Figure~\ref{fig:4S2}). At a coarse-grained level, this structure is analogous to a TSSA, with each super-node representing the collective of genes defining one phenotype. To systematically study how regulatory noise interacts with this architecture, we generated ensembles of artificial networks using the stochastic block model~\cite{Holland1983STOCHASTIC}, constructing networks in which intra-team edges are activating and inter-team edges are inhibitory (Figure~\ref{fig:5}A, leftmost panel). We characterised each network by two parameters: its edge density, defined as the fraction of possible edges that are present, and its impurity, defined as the fraction of edges that conflict with the expected team structure (i.e., activating edges between teams or inhibiting edges within teams). At an impurity of 0.5, the network becomes equivalent to an Erd\H{o}s--R\'{e}nyi random network with equal probability of activating and inhibiting edges. Biologically occurring binary cell-fate networks have densities $\leq 0.3$ and impurities of approximately 5\% (EMT network: density 0.23, impurity $\sim$0.05; Gonadal network: density 0.17, impurity $\sim$0.05), providing natural reference points for the artificial network parameter space. We generated 50 networks per condition and simulated each under the three noise types at the noise levels described previously. Given the size of these networks, exploration of regulatory noise using similar ODE based models as with small-scale networks earlier would be computationally infeasible. Hence, we utilized a threshold based boolean formulation to analyze these networks instead \cite{FontClos2018}. In this formalism, each node is allowed to take either -1 (low) or +1 (high) expression levels. The edges of the network are given an weight of +1 (activation) or -1 (inhibition). State of the network is updated asynchronously, based on an ising-model inspired, additive rule (see Methods). The phenotypic state of each network, state vectors were converted into team scores --- the fraction of nodes in each phenotypic team that are active --- and a phenotypic score defined as the difference between the two team scores. The stochastic block model pre-determines the team identity of each node. Values near $\pm 1$ indicate terminal phenotypes (epithelial/mesenchymal or Sertoli/Granulosa), while values near 0 indicate hybrid states.

% \begin{figure}
%     \centering
%     \includegraphics[width=\linewidth]{Figs/Figure5.jpg}
%     \caption{\textbf{Analysis of large-scale binary cell-fate decision networks.} A) Schematic of the block-wise creation of artificial teamed networks with impurity, with ER random network for comparision at the right end. B-D) Distributions of phenotypic score across 50 networks X 100 trajectories, for B) Additive noise with no impurity, C) Additive noise with varying impurity at a constant density of 0.1 and D) Multiplicative noise at a constant density of 0.1. The color represents the noise level, ranging from 0 to 0.05. }
%     \label{fig:5}
% \end{figure}

% \paragraph{Effect of density.}
% At zero impurity, the phenotypic score distributions under additive noise show a clear density dependence (Figure~\ref{fig:5}B). At low density ($\rho = 0.1$), noise levels of $\sigma \geq 0.005$ are sufficient to drive the system almost entirely into hybrid states, with the phenotypic score distribution collapsing toward zero. At higher densities ($\rho = 0.3$ and $0.5$), terminal state mass is progressively better preserved, though additive noise at $\sigma = 0.05$ still produces substantial hybrid state accumulation even at $\rho = 0.5$. Density therefore delays but does not prevent the transition to hybrid-dominated dynamics under additive noise. Multiplicative and stationary noise, by contrast, maintain terminal states robustly across the full density range explored, even at the highest noise levels tested (Figures~\ref{fig:5S1}A(iii)--(iv), \ref{fig:5S2}A).

% \paragraph{Effect of impurity.}
% Even in the absence of noise, increasing impurity progressively erodes the terminal phenotype frequency, with terminal states lost entirely at impurity 0.3 for networks with density 0.1 (Figure~\ref{fig:5}C). Under additive noise, this erosion is accelerated: at density 0.1, noise $\sigma \geq 0.005$ drives the system fully into hybrid states at impurity 0, and the additional effect of increasing impurity is to flatten the residual polarization of the phenotypic score distribution toward $\pm 0.5$--$0.75$, producing a distribution tightly localised near zero (Figures~\ref{fig:5}C, \ref{fig:5S1}A(ii)). Multiplicative noise again shows a qualitatively different behaviour: terminal states are maintained across the full range of impurity levels tested, even at $\sigma = 0.05$ (Figure~\ref{fig:5}D). Stationary noise behaves similarly to multiplicative noise, preserving terminal states robustly across both density and impurity ranges (Figure~\ref{fig:5S2}). The 2D team score distributions (Figure~\ref{fig:5S1}) provide the clearest summary of this contrast: under additive noise at $\sigma = 0.05$, all conditions converge to a central $(0.5, 0.5)$ mass in team score space regardless of density or impurity, while multiplicative and stationary noise maintain corner-localised distributions corresponding to terminal phenotypes.

% Together, these results indicate that the noise-type asymmetry identified in small motifs --- additive noise driving hybrid state accumulation, multiplicative and stationary noise preserving terminal states --- is a general property of teamed network architectures rather than a feature specific to the toggle switch or TSSA topology. Network density and impurity modulate the threshold at which additive noise produces this effect, with sparser and more impure networks being more susceptible, but neither parameter rescues terminal state stability under strong additive noise, nor does either parameter destabilise terminal states under multiplicative or stationary noise. We note that in the absence of continuous noise, static reductions in edge weights also shift phenotypic distributions toward hybrid states, though the effect is less pronounced than under continuous noise; this analysis is presented in full in the Supplementary Material (Section~\ref{ssec:stat_pert}).

% \subsubsection{Biological networks}

Large-scale GRNs underlying binary cell-fate decisions share a common coarse-grained structure: nodes can be partitioned into two ``teams'', with intra-team interactions predominantly activating and inter-team interactions predominantly inhibitory (Figure~\ref{fig:5S1}). At a coarse-grained level, this structure is analogous to a TSSA, with each super-node representing the collective of genes defining one phenotype. The biological GRNs adhere to this structure at a coarse-grained level as well. The EMT network, a 23-node network underlying epithelial--mesenchymal transition (EMT)~\cite{Huang2017}, has two teams of nodes corresponding to the epithelial and mesenchymal phenotypes respectively. In fact, we previously showed that all large-scale EMT GRNs constructed previously could be coarse-grained in two mutually inhibiting teams, and that the presence of these teams gives rise to a robust, low-dimensional phenotypic landscape despite the size of these networks~\cite{hari2022landscape, Hari2025LowNetworks}. We also showed that these teamed networks are robust to noise in the network state, and extensive mutations disrupting the network of positive feedback present in these GRNs are required to stabilize what is otherwise weakly stable hybrid phenotype~\cite{rashid2022network, Hari2025APlasticity}. The 19-node gonadal cell-fate decision network cell-fate decisions between Sertoli and Granulosa fates that are upstream of male and female development respectively~\cite{rios2015boolean} has two teams of nodes corresponding to Sertoli and Granulosa phenotypes, and carry the advantages of team structure across simulation formalisms~\cite{Hari2025LowNetworks, Shyam2025MutuallyNetworks}.

Given the size of these networks, exploration of regulatory noise using similar ODE based models as with small-scale networks earlier would be computationally infeasible. Hence, we utilized a threshold based boolean formulation to analyze these networks instead \cite{FontClos2018}. In this formalism, each node is allowed to take either -1 (low) or +1 (high) expression levels. The edges of the network are given a weight of +1 (activation) or -1 (inhibition). The state of the network is updated asynchronously, based on an ising-model inspired, additive rule (see Methods). The binary state vectors ($S \in \{0, 1\}^N$) were converted into team scores --- the fraction of nodes in each team that are active --- and a phenotypic score defined as the difference between the two team scores. Values near $\pm 1$ indicate terminal phenotypes (epithelial/mesenchymal or Sertoli/Granulosa), while values near 0 indicate hybrid states. Under deterministic conditions, the phenotypic landscape is largely concentrated in regions corresponding to terminal states with team scores $\pm$ 1. 

\begin{figure}
    \centering
    \includegraphics[width=\linewidth]{Figs/Figure4.jpg}
    \caption{\textbf{Regulatory noise in large-network boolean models} A) Representative trajectories of the phenotypic score for EMT network (i) and Gonadal network (ii) at additive noise of 0.05 magnitude. B) Density of team scores across the representative trajectory of the EMT network, with frequency represented on a scale from low (blue) to high (yellow). WT represents the frequency of team score configurations obtained from the steady states of the unperturbed network, while Rand represents that of the simulations corresponding to the trajectories in A. C) Relative frequency (equivalent to MRT) of different phenotypic scores across 1000 simulation trajectories of additive noise over simulation time for (i) EMT network, (ii) Gonadal network. D) Same as C, but for Multiplicative and Stationary noise for EMT network.}
    \label{fig:5}
\end{figure}

Under additive noise, both biological networks show behaviour consistent with the small network analysis. Trajectories initiated from terminal phenotypes of the unperturbed network rapidly migrate to hybrid phenotypic scores and remain there throughout the simulation, regardless of initial condition (Figure~\ref{fig:5}A). The transition occurs sharply: even at $\sigma = 0.001$, terminal state frequencies are markedly reduced (Figure~\ref{fig:5}B), and by $\sigma = 0.005$ the phenotypic score distribution has saturated into a stable hybrid-dominated configuration that changes little with further increases in noise (Figure~\ref{fig:5}C, \ref{fig:5S2}A). Multiplicative and stationary noise produce minimal disruption to the phenotypic landscape, with trajectories retaining high terminal state frequency similar to the unperturbed network across all noise levels tested (Figure~~\ref{fig:5}C, \ref{fig:5S2}B, C).

The two networks differ, however, in the characteristics of their hybrid states. In the EMT network, the hybrid phenotypic score distribution retains some polarisation toward $\pm 0.5$--$0.75$, reflecting a tendency for one team to remain more active than the other even in the hybrid regime (Figure~\ref{fig:5}C(i), B). In the Gonadal network, the distribution collapses tightly to near zero, indicating near-equal team expression with no residual polarization (Figure~\ref{fig:5}C(ii)). Both networks have similar densities and team structures to each other,  suggesting that network size and specific wiring topology beyond the team structure are likely contributors to this difference. A full characterization of the structural features that determine the degree of polarization within the hybrid regime remains an open question. Overall, our results indicate that regulatory noise can have a significant impact on the phenotypic landscape, with the direction of the impact determined by the nature and the magnitude of the noise. While additive noise can drive the cell-fate decision system up towards an undifferentiated phenotype, multiplicative noise can drive the system down towards a terminally differentiated state (Figure~\ref{fig:schematic}). 
\begin{figure}
    \centering
    \includegraphics[width=\linewidth]{Figs/Schematic.png}
    \caption{\textbf{Graphical representation of the effect of regulatory noise on the cell-fate decision landscape supported by GRNs}}
    \label{fig:schematic}
\end{figure}

\section{Discussion}

Regulatory interactions in gene networks are rarely static. Transcription factor binding affinities fluctuate due to conformational heterogeneity, post-translational modification, and competition for genomic binding sites, yet GRN models overwhelmingly treat these interaction strengths as fixed parameters. Static models have nonetheless been remarkably successful at predicting cell-fate outcomes, from EMT circuits to Boolean perturbation frameworks. This success is not a reason to dismiss stochasticity in kinetic parameters. Instead, we must ask what such noise would be mistaken for if it were present. Here, we studied the consequences of allowing the fold-change parameter which sets the regulatory efficacy of each interaction to vary stochastically over time, across cell-fate decision motifs of increasing complexity and two well-characterised biological GRNs (EMT and gonadal cell-fate determination).

\paragraph{Regulatory noise can cloud the inference of network topology}
Prior theoretical work has shown that noise can shift or stretch the effective bistable range of a network relative to its deterministic prediction, and can generate bimodal, switch-like population distributions in circuits that are deterministically monostable for all parameter values~\cite{OchabMarcinek2010w}. Regulatory noise of the kind we studied here also modifies the qualitative details of the phenotypic landscape. The (1,1) co-expression state, which is absent or negligible in the steady state distribution emergent from toggle switch in the deterministic limit, becomes dominant under additive noise, with mean residence time (MRT) rising to $\sim$0.6 for monostable parameter sets and $\sim$0.85--0.9 for bistable ones by $\sigma = 0.1$. This suggests that a positive correlation between two genes observed in a population snapshot like transcriptomic data would not be able to distinguish between mutual activation under deterministic conditions and mutual inhibition under stochastic conditions. Similarly, the single-high states that are not stable in larger networks such as the toggle tetrahedron become dominant under multiplicative noise. Previous simulations have shown that co-existence of single-high states beyond three node networks is not feasible. A network's population-level phenotype distribution can therefore reflect its noise regime as much as its wiring diagram. Inferring topology from steady-state expression data without accounting for regulatory noise therefore risks systematically over- or under-estimating regulatory complexity --- precisely the motivation for modelling regulatory noise explicitly rather than treating it as a second-order correction.

\paragraph{The three noise modes map onto different modes of action of IDP transcription factors}
The three noise modes studied here - stationary, additive and multiplicative each corresponds to a mechanistically distinct route by which an intrinsically disordered transcription factor's conformational ensemble can modulate its effective regulatory strength. Each noise mode redistributes state occupancy in a qualitatively different direction.

Stationary noise, where $\lambda$ fluctuates around an initially sampled $\lambda$ value, has the smallest effect of the three modes at low-to-moderate $\sigma$. But it is not completely inert: at the highest noise levels tested, it produces a modest rise in hybrid-state MRT that is directionally consistent with, though substantially weaker than, the additive noise effect. This is the same boundary-driven mechanism as additive noise, operating from a much higher noise threshold rather than being a qualitatively distinct, noise-immune regime. Molecularly, this mode corresponds to the inherent disorderedness present in TFs, which does not affect the TF's binding to the promoter, but is theorized to help flexibly search for and recruit the transcriptional machinery~\cite{Holehouse2023TheRegions, Jonas2025IntrinsicallySearch}. The mean $\lambda$ that this distribution reverts to represents a deep valley in the TF-promoter-transcription machinery ensemble, with the noise representing minor fluctuations around the valley that can be captured by a deterministic description of the system.  

Multiplicative noise, where perturbations scale with the current value of $\lambda$ such that inhibitory $\lambda$ is driven toward zero, single-high states are reinforced or actively restored (even the rare deterministic TS models that permit (1,1) show that state suppressed), and in TT, double-high hybrid states present in the deterministic attractor are converted to single-high states. Molecularly, this models a higher disorderedness in TFs than represented by Stationary noise, and captures the conformational sampling that progressively locks the ensemble toward one sub-state. While this behavior is not well studied in the context of TFs, specific post-translational modifications, particularly phosphorylation, stabilise particular IDP conformers with half-lives of minutes to hours~\cite{Kulkarni2022, Kulkarni2019}. Because each such modification event can make subsequent modification more likely, the process compounds in a state-dependent, self-reinforcing manner --- the same qualitative signature multiplicative noise produces on $\lambda$. 

Additive noise, in which $\lambda$ performs a bounded random walk, evolving into a distribution independent of its initial value. Additive noise progressively weakens inhibitory interactions and destabilises single-high attractors in favour of hybrid co-expression states, an effect that intensifies with network complexity, with the MRT of all-high states reaching ~1.0 at high $\sigma$ as long as the state is accessible by the parameter set. Molecularly, this maps most directly onto unconstrained conformational diffusion agnostic to any contact with the DNA: when an IDP's disordered region samples its ensemble without being biased toward any particular sub-state, its influence on effective fold-change varies in a manner unanchored to the interaction's current value --- precisely the additive case~\cite{Ferrie2022Structure-functionRegulation, Holehouse2023TheRegions}. 


Together, these three molecular readings recover the diversity of IDP-mediated fold-change fluctuation --- unbiased, locking, and state-anchored --- that motivated our use of multiple noise modes rather than a single generic noise term.


\paragraph{The stochastic-vs-deterministic comparison indicates a genuine dynamical contribution, not merely a redistribution of parameters.}
A natural objection to the additive and multiplicative noise results above is that they might simply reflect resampling: if additive noise broadens or shifts the effective distribution of $\lambda$ values explored, then running deterministic RACIPE on parameter sets drawn from that same noise-induced $\lambda$ distribution might reproduce the same state occupancies without requiring any dynamical argument. We tested this directly by comparing the MRT obtained from actual stochastic trajectories against the MRT obtained from deterministic simulations sampled from the stationary $\lambda$-distribution that each noise mode generates. 
%For multiplicative noise, and for fluctuating noise at low-to-moderate $\sigma$, the two approaches agree closely, indicating that in these regimes the system is close to adiabatic: expression tracks its instantaneous equilibrium at each sampled $\lambda$, and the outcome is fully explained by where $\lambda$ spends its time. 
For both additive and multiplicative noise, the stochastic MRT of high-expression (hybrid) states and single-high (terminally differentiated) states respectively, exceeds what the $\lambda$-distribution alone predicts. This indicates that additive and multiplicative noise contribute a dynamical effect beyond parameter redistribution --- consistent with the system being pushed away from the adiabatic limit, so that the temporal structure of the $\lambda$ trajectory itself, not merely its stationary distribution, shapes which states are visited. This distinction matters for interpretation: it means additive and multiplicative regulatory noise are not simply equivalent to sampling from a noisier RACIPE parameter ensemble, and that non-adiabatic coupling between conformational and expression timescales is itself a biologically relevant variable.

\paragraph{The effect generalises beyond small motifs to two biological GRNs.}
To test whether these motif-level results are a generic property of larger biological networks rather than an artefact of the toggle switch or triad specifically, we simulated two real biological GRNs: a 23-node EMT network and a 19-node gonadal cell-fate network (Sertoli/Granulosa determination), using a Boolean formalism, under the same three noise modes. Both showed the same qualitative behaviour as the small network motifs, and the transition was sharp. Under additive noise, terminal-state frequency was already markedly reduced at $\sigma = 0.001$ and the phenotypic score distribution had saturated into a hybrid-dominated configuration by $\sigma = 0.005$, essentially independent of initial condition. Multiplicative and stationary noise again produced minimal disruption. The two networks differed, however, in the character of their hybrid states: the EMT network retained partial polarisation ($\pm 0.5$--0.75) even in the hybrid regime, indicating one team tends to remain more active than the other, whereas the gonadal network collapsed fully to the (0.5, 0.5) team-score centre. Both networks have comparable density and team strucutres, indicating that this difference must be the property of other network properties, such as network size or specific wiring beyond team structure, that we have not yet characterized (see Limitations).

\paragraph{Hybrid states, entropy, and pathological implications.}
A unifying consequence of additive regulatory noise, across every topology and scale studied here, is an increase in phenotypic entropy: the system spends progressively more time in co-expression states that are absent from, or negligible in, the deterministic attractor landscape. In the EMT network specifically, these co-expression states correspond directly to the hybrid epithelial/mesenchymal (hybrid E/M) phenotype. The hybrid E/M state is now well-established as a functionally distinct and clinically relevant cell state: hybrid E/M cells display elevated stemness, enhanced collective migration, greater metastatic competence, and stronger therapy resistance relative to cells occupying either terminal phenotype~\cite{Pastushenko2019EMTMetastasis, jolly2015implications, pasani2020hybrid, Sahoo2021}. Our results suggest a concrete biophysical route through which hybrid E/M incidence could be elevated in tumours: if regulatory noise in regulatory interactions is increased --- by the tumour microenvironment, or by non-genetic or genetic alterations to TF binding --- additive-like fold-change fluctuations are, on their own, sufficient to drive both an isolated bistable motif and a full 23-node biological network into a predominantly hybrid regime, with a transition threshold ($\sigma \sim 0.001$--0.005) well within a biologically plausible range. This directly links molecular-scale disorder in transcriptional regulation to population-level phenotypic heterogeneity of the kind that fuels tumour progression and drug tolerance.

\paragraph{Stem cell identity and chromatin state.}
There is a strong correspondence between our model and developmental biology that extends to pluripotency. We build on the established view that hybrid, co-expressing states reflect a less differentiated, higher-entropy configuration, while single-high terminal states reflect the low-dimensional, more homogeneous endpoint of differentiation~\cite{Teschendorff2017Single-cellTranscriptome, Hari2025LowNetworks}. Pluripotent stem cells maintain a globally open, decondensed chromatin state supporting stochastic, broadly accessible transcription~\cite{Meshorer2006HyperdynamicCells, Efroni2008GlobalCells}, consistent with recent evidence that fluctuations in chromatin accessibility are actively amplified to permit ectopic, multi-lineage-like gene expression programs in a physiological setting~\cite{Gamble2025}. In our framework, this corresponds to an additive-noise regime: open chromatin raises the density of accessible genomic sites, increasing the likelihood that a TF engages non-specific sites and modulates its effective regulatory strength in a manner unanchored to its current interaction state.

Our results yield a distinct, testable prediction from this correspondence: transition towards a terminally differentiated state does not require a specific instructive signal to execute the transition itself. Instead, we propose a two-component picture : extrinsic signals bias which hybrid or lineage-primed state a progenitor occupies, while a shift toward a multiplicative-noise-like regime supplies the drift that resolves that hybrid state into a single-high terminal fate, without requiring the signal to specify the differentiation step itself. This separates two components of differentiation that are typically treated as a single instructive event: fate selection (signal-dependent) and fate resolution (noise-mode-dependent, signal-agnostic). This view helps reconcile the longstanding debate over whether lineage commitment is fundamentally instructive or noise-driven~\cite{<lineage commitment uncertainty refs>}, and offers a mechanistic account of how both could be simultaneously true at different stages of the same transition.

\paragraph{IDPs, conformational noise, and cancer.}
Cancer mutations targeting IDPs are of particular relevance given the mapping above. Mutations within disordered regions can increase non-specific binding affinity by expanding the conformationally accessible ensemble, and can introduce new disordered segments into previously structured TFs, collectively amplifying the magnitude of additive-like fold-change fluctuations~\cite{Kulkarni2023IntrinsicallyLater}. Under our framework, such mutations are predicted to increase hybrid-phenotype frequency and phenotypic heterogeneity within a tumour even without any change to deterministic network topology --- a functional interpretation of IDR mutations beyond direct loss-of-function, and a motivation for quantifying the conformational-noise contribution of specific cancer-associated IDR variants.

\paragraph{Decoupling conformational drift from binding: justifying the bounds.}
Our model decouples the noise in regulatory strength from the binding event itself. From the IDP perspective, conformational drift need not be gated by whether the transcription factor is currently bound; the disordered region can sample its accessible ensemble whether or not it is engaged at a promoter. What is bounded is not the conformational search but the downstream consequence of any given conformation: there is an upper limit on how efficiently any TF--DNA interaction can drive transcription, set by factors external to the TF's own conformational landscape, such as polymerase recruitment kinetics and the physical throughput of the transcription machinery, and symmetrically a floor, since a fold-change cannot fall below the point of no functional engagement. These bounds are therefore a direct outcome of treating conformational sampling and transcriptional output as separate processes linked only through a saturating map, not a numerical convenience. This is also why we implement the bounds as hard limits rather than a soft, binding-aware boundary: a boundary that softens as it is approached would require the underlying mechanism to sense the binding event it is nominally decoupled from, which we consider a less general assumption than a hard ceiling on function that is agnostic to how it is reached.

\paragraph{The sign-change assumption.}
A consequence of this implementation is that inhibitory fold-changes ($\lambda \in [\lambda_{min},1]$) and activatory fold-changes ($\lambda \in [1, \lambda_{max}]$) occupy disjoint intervals under every noise mode studied here: no amount of additive, multiplicative, or stationary fluctuation can carry an interaction across $\lambda$=1 to invert its sign. Continuous modulation of an existing interaction's strength (this study) and discrete topological rewiring, including sign inversion (deferred), are therefore disjoint mechanisms by construction, not merely different in degree. A molecular route to genuine sign inversion is nonetheless well documented: transcription factors such as RBPJ switch between repressive and activating output depending on which corepressor or coactivator complex they currently engage~\cite{Dong2026MolecularFunctions}, a behaviour formalised as "incoherent" regulation, in which a single factor acts as both activator and repressor at a given target under a fixed molecular context~\cite{Martinez-Corral2025EmergenceContext}. Because our additive-noise results show hybrid states increasing in frequency without any edge changing sign, we read this as continuous noise \emph{phenocopying} part of the phenotypic signature normally attributed to rewiring, not as evidence that rewiring has occurred. Extending the present formalism to allow genuine sign inversion, and quantifying how much of the rewiring phenotype can be reproduced by noise alone, is left to future work.

\paragraph{Limitations and future directions.}
Several aspects of this study warrant further investigation. First, the fold-change noise model is, by design, abstract in its mechanistic scope: it captures the effect of molecular fluctuation on regulatory strength without specifying its precise source, which limits direct quantitative comparison with specific experimental measurements of TF binding dynamics; mapping our noise modes onto measurable quantities such as single-molecule TF residence times or IDR conformational correlation times would strengthen this connection. Second, our formalism perturbs fold-change parameters, which is distinct from classical state-space (expression-level) noise of the Langevin type; because such noise acts within a fixed landscape and is continuously damped back toward the current attractor. We expect it to have a comparatively limited effect relative to fold-change noise, which can deform the landscape itself. Studies of noise-driven differentiation predict a rescaling of the existing attractors in the landscape~\cite{Herrera2024ADifferentiation, Mojtahedi2016CellTransition}. We had previously shown that combining state-space noise with edge perturbations can introduce new dynamical behavior in the same attractor landscape \cite{Hari2025APlasticity}, which is closer to our predictions here, but fundamentally different still, as the fold-change perturbations generate new attractors in the landscape (hybrid states in additive noise and terminal states in TT4 in multiplicative noise). A formal, sensitivity-matched comparison between the two noise classes remains future work. Third, the treatment of activatory interactions under multiplicative noise requires further development, since the asymmetry between inhibitory and activatory parameters under state-dependent noise introduces complexities the current formalism does not fully address. Fourth, for the large biological networks our analysis is limited to the Boolean framework, which sacrifices quantitative resolution of expression levels; an ODE-based treatment under parameter noise, while computationally demanding, would allow more direct comparison with the small-motif results and with single-cell expression data. Fifth, the difference in hybrid-state polarization between the EMT and gonadal networks under additive noise remains mechanistically unexplained and points to structural features beyond density and impurity that we have not yet characterised. Sixth, the current formalism applies noise uniformly across all regulatory edges; because the aberrant conformational dynamics associated with IDP mutation are, mechanistically, localised events, only a subset of interactions in a large network may be affected at any given time. Empirical estimates of interactor turnover under a single phosphorylation-state transition in PAGE4/MED25-driven regulation suggest that roughly a quarter of stable interactions can be affected by a single such event, which we take as a rough guide to the scale of localisation worth exploring in future work, ideally validated on small networks before being extended to the large biological networks studied here. Finally, our present results establish the direction and generality of each noise mode's effect but do not yet resolve how noise mode interacts with multistability itself --- for instance, whether multiplicative noise can convert monostable parameter regimes into multistable ones by resolving symmetric attractors stochastically, or whether this generalises across the full network set we have simulated; these results require further verification and are not reported here.

\section{Methods}

\subsection{Network models and RACIPE ensemble generation}\label{sec:RACIPE}

Each network motif is modelled as a system of ordinary differential equations (ODE), one ODE per node, following the shifted Hill formalism of RACIPE~\cite{Huang2017}:
\begin{equation}
\frac{dX}{dt} = g_X \prod_{i \in P} \frac{H(Y_i, Y_{i0}, n_{Y_i}, \lambda_{Y_i})}{\lambda_{Y_i}} \prod_{i \in N} H(Y_i, Y_{i0}, n_{Y_i}, \lambda_{Y_i}) - d_X X
\label{eq:odeM}
\end{equation}
\begin{equation}
H(Y_i, Y_{i0}, n_{Y_i}, \lambda_{Y_i}) = \lambda_{Y_i} \frac{Y_i^{\,n_{Y_i}}}{Y_{i0}^{\,n_{Y_i}} + Y_i^{\,n_{Y_i}}} + \frac{Y_{i0}^{\,n_{Y_i}}}{Y_{i0}^{\,n_{Y_i}} + Y_i^{\,n_{Y_i}}}
\label{eq:hillM}
\end{equation}
where $g_X$ and $d_X$ are the production and degradation rates of node $X$, $P$ and $N$ index its activating and inhibiting regulators, and $\lambda_{Y_i}$ is the fold-change by which regulator $Y_i$ modulates $X$'s production ($\lambda_{Y_i} > 1$ for activation, $\lambda_{Y_i} < 1$ for inhibition).

For each network topology, RACIPE~\cite{Huang2017} was used to generate an ensemble of $10{,}000$ parameter sets by uniform sampling within standard biologically plausible ranges, with each parameter set solved from 100 random initial conditions to identify its steady-state attractor(s). The resultant ensemble of steady states was used to calculate thresholds of individual node expressions ($\theta_X$) as the average expression across the ensemble, which were then used to discretize the steady states into strings of 1's and 0's. Based on the number of such discretized steady states allowed, each parameter set was assigned a multistability class (monostable, bistable, tristable, tetrastable). A representative subset of 100 parameter sets was selected from each stability class present for a given network for downstream stochastic analysis. 

\subsection{Regulatory noise model} \label{sec:noise}

To simulate regulatory noise, we allowed the fold-change parameters $\lambda_{Y_i}$ to vary stochastically over time rather than remaining fixed at their RACIPE-sampled value, using a periodic update applied every $dt = 0.01$ (a.u.) during ODE integration. Three noise modes were implemented, differing in how each update depends on the parameter's current value:
\begin{align}
\text{Additive:} \quad & \lambda(t+dt) = \lambda(t) + \mathcal{N}(0, \sigma_{\text{eff}}^2), \\
\text{Multiplicative:} \quad & \lambda(t+dt) = \lambda(t)\left(1 + \mathcal{N}(0, \sigma_{\text{eff}}^2)\right), \\
\text{Stationary:} \quad & \lambda(t+dt) = \lambda(0) + \mathcal{N}(0, \sigma_{\text{eff}}^2),
\end{align}
where $\sigma_{\text{eff}} = \sigma \cdot \lambda_{\max}$, with $\lambda_{\max}$ the upper bound of the relevant fold-change (inhibitory: $\lambda_{\max} = 1$, so $\sigma_{\text{eff}} \approx \sigma$; activatory: $\lambda_{\max} = 100$), so that a given nominal noise level $\sigma$ corresponds to a comparable proportion of each parameter's admissible range regardless of whether it governs activation or inhibition. After each update, $\lambda$ was clamped to its biologically admissible bounds ($[0.001, 0.999]$ for inhibitory, $[1.001, \lambda_{\max}]$ for activatory parameters). Noise levels $\sigma \in \{0, 0.001, 0.005, 0.01, 0.05, 0.1, 0.5, 1\}$ were tested for each noise mode and network.

\subsection{Stochastic simulation and mean residence time}

For each selected parameter set and noise level, 25 independent stochastic trajectories were simulated, each initialised from an independently drawn initial condition $u_i(0) = U(0,1) \times (g_i/d_i)$ for every node $i$ (i.e.\ uniformly between zero and that node's maximal attainable expression). Trajectories were integrated over $t \in [0, 1000]$ (a.u.) using the Tsit5 adaptive-step solver (relative tolerance $10^{-4}$, absolute tolerance $10^{-5}$). Each trajectory was discretized into a sequence of Boolean node states using the RACIPE-derived thresholds, and the first 30\% of each trajectory was discarded as transient. Mean residence time (MRT) for a given discrete state was computed as the fraction of the remaining trajectory (across all 25 stochastic realizations) spent in that state.

\subsection{Deterministic $\lambda$-resampling comparison}

To test whether a given noise mode's effect on state occupancy could be explained purely by the stationary distribution of $\lambda$ it induces, rather than by any dynamical (non-adiabatic) contribution, we compared stochastic MRT against a deterministic control: for each parameter set and noise level, $\lambda$ values were resampled from the empirical stationary distribution generated by that noise mode (via \texttt{get\_lambda\_dist}), and the corresponding deterministic steady state was solved; this was repeated over $\text{DET\_ITERS} = 10$ resampled draws per parameter set, and the resulting basin-occupancy fractions (MRT$_{\text{det}}$) were compared directly against the stochastic MRT from the same parameter set and noise level.


\subsection{Boolean simulations of large cell-fate networks}

For larger networks, we used Boolean modeling to simulate network dynamics, as opposed to the ODE models of RACIPE. We used a threshold-based Boolean formalism \cite{FontClos2018}, inspired by the Ising model of spin-glass systems. A network is defined by $N$ nodes, $E$ edges, and a topology, i.e., the arrangement of $E$ edges among $N$ nodes. The Ising model employs discretized expressions of nodes, with $s_i \in \{-1, 1\}$; $i \in 1, \ldots, N$, where $s_i = 1$ and $s_i = -1$ represent an active and inactive node, respectively. A state at any time $t$ is the set of expression levels of all nodes, $S(t) = \{s_i(t)\}$; $i \in 1, \ldots, N$.

The system state is updated in discrete time steps using an asymmetric adjacency matrix ($\adj$), which stores the nature of pairwise interactions between nodes. In a static network, activating and inhibiting edges carry edge weights of $+1$ or $-1$, respectively, with $0$ indicating no interaction. The expression level of node $i$ at time $t$ is given by:
\begin{equation} \label{eq:bool}
    s_i(t) = 
\begin{cases}
    +1, & \text{if } \sum_j \adj_{ij} s_j(t-1) > 0 \\
    -1, & \text{if } \sum_j \adj_{ij} s_j(t-1) < 0 \\
    s_i(t-1), & \text{if } \sum_j \adj_{ij} s_j(t-1) = 0
\end{cases}
\end{equation}
where $\sum_j \adj_{ij} s_j$ represents the net regulatory input to node $i$. A steady state is defined as a state $S(t)$ satisfying $S(t+1) = S(t)$.

\subsection{Large network metrics}

\subsubsection{Influence matrix}
The influence matrix is formulated as:
\begin{equation}
    \inm_{l_\mathrm{max}} \equiv \inm = \frac{\sum_{l=1}^{l_\mathrm{max}} \frac{\adj^l}{\adj_{\mathrm{max}}^l}}{l_\mathrm{max}},
\end{equation}
where $l$ is the path length, representing a path of $l$ edges connecting a pair of nodes, and $l_\mathrm{max} = 10$ is the maximum path length considered. The influence matrix captures the net bidirectional effect between a pair of nodes mediated by all paths of length $l \in \{1, \ldots, l_\mathrm{max}\}$. $\adj_{\mathrm{max}}^l$ is derived from $\adj^l$ by setting all non-zero entries to $1$, and the division $\frac{\adj^l}{\adj_{\mathrm{max}}^l}$ is element-wise.

\subsubsection{Teams}
The wild-type influence matrices of the biological networks considered in this study are displayed in Figure~\ref{fig:5S1}. Each network can be partitioned into two ``teams'' of nodes, defined such that the influence between any pair of nodes is positive if they belong to the same team and negative otherwise. The team strength of a two-teamed network is defined as:
\begin{align} \label{eq:ts_def}
    \ts_{kl} &= \frac{\sum_{i \in T_k, j \in T_l} \inm_{ij}}{n_{kl}}, \quad k, l \in \{1, 2\} \\
    \ts &= \frac{\sum_{k,l \in \{1,2\}} |\ts_{kl}|}{4},
\end{align}
where $T_k$ is the $k$-th team.

\noindent\textbf{\textit{Team expression}}\label{text:team_exp}: For a given state $S$, team expression is defined as the fraction of nodes belonging to a team that are active (i.e., $s_i = 1$).

% \paragraph{Artificial networks} For artificial networks, the team identities are decided before generating the networks, and they are used directly without requiring the calculation of influence matrix.

\noindent\textbf{\textit{Nature of stable states}}: Stable states of two-team networks are classified as either dominant (most frequent) or hybrid (non-dominant). Due to the bimodal phenotypic distribution characteristic of toggle switch-like networks, dominant states are those in which only one team is active, i.e., states with team expressions of $(0,1)$ and $(1,0)$.

\subsection{Stochastic simulation of large networks}

\subsubsection{Inducing noise}\label{sssec:noise_def}

In biological networks, the strengths of regulatory interactions are subject to temporal variation arising from fluctuations in protein concentrations, post-translational modifications, and other molecular sources. We model this regulatory noise by perturbing the edge weights of the adjacency matrix over time. Three noise modes are implemented, mirroring those used in the ODE-based simulations of small motifs (Section~\ref{sec:RACIPE}):

\paragraph{Additive (state-independent) noise}
At each time step, a perturbation drawn from $\mathcal{N}(0, \sigma^2)$ is added to each edge weight independently of its current value:
\begin{equation}
    \mathrm{EW}(t) = \mathrm{EW}(t-1) + \delta, \quad \delta \sim \mathcal{N}(0, \sigma^2).
\end{equation}
Edge weights are clipped to $[-1, 1]$ after each update to preserve the biological interpretation of interaction strengths.

\paragraph{Multiplicative (state-dependent) noise}
At each time step, each edge weight is scaled by a factor drawn from $\mathcal{N}(1, \sigma^2)$, so that the magnitude of the perturbation scales with the current edge weight:
\begin{equation}
    \mathrm{EW}(t) = \mathrm{EW}(t-1) \cdot (1 + \delta), \quad \delta \sim \mathcal{N}(0, \sigma^2).
\end{equation}

\paragraph{Stationary (mean-reverting) noise}
At each time step, each edge weight is redrawn as its wild-type value plus an independent perturbation, so that the weight fluctuates around its nominal value without drift:
\begin{equation}
    \mathrm{EW}(t) = \mathrm{EW}_0 + \delta, \quad \delta \sim \mathcal{N}(0, \sigma^2),
\end{equation}
where $\mathrm{EW}_0$ denotes the wild-type edge weight. Unlike the additive mode, this formulation prevents long-range drift in parameter values and models a situation where molecular fluctuations are transient and the nominal interaction strength is effectively restored between perturbation events.

In all three modes, $\sigma^2 = 0.01$ is used as the noise variance. Each biological network is simulated for $1000$ trajectories of $10{,}000$ time steps each; artificial networks are simulated for $100$ trajectories of $1{,}000$ time steps each.

\subsubsection{Method of state update}
In practice, the system can be updated in two ways \cite{garg2008synchronous, schwab2020concepts}:
\begin{itemize}
    \item \textbf{Synchronous update}: At each time step, every node is updated simultaneously using Eq.~\ref{eq:bool}.
    \item \textbf{Asynchronous update}: At each time step, a single randomly chosen node is updated using Eq.~\ref{eq:bool}.
\end{itemize}
While synchronous updates are computationally efficient, they can produce limit cycles from which the system cannot escape. Asynchronous updates are less efficient but avoid limit cycles and implicitly capture reaction noise by assuming that only one regulatory event occurs per time step. To combine the benefits of both schemes, we check for steady states synchronously every $200$ time steps while updating the network asynchronously. For continuous perturbation simulations, synchronous updates are used throughout, as persistent perturbations make entrapment in limit cycles unlikely.

\subsubsection{Initial conditions}

Peripheral nodes — those involved in interactions without feedback — are excluded from the analysis, as the focus of this study is on teaming behaviour arising from mutual regulatory effects. Initial conditions are drawn from the terminal states and the most dominant hybrid states identified in the unperturbed network's steady-state distribution.

\section{Acknowledgments}

MKJ was supported by the Ramanujan Fellowship (SB/S2/RJN-049/2018) awarded by the Science and Engineering Research Board (SERB), Department of Science and Technology, Government of India. MKJ was also supported by Param Hansa Philanthropies. KH and HL were supported by the Center for Theoretical Biological Physics, NSF PHY-2019745, and under Award Number MCB-2114191

\printbibliography

\beginsupplement

\section{Derivation of the shifted Hill functions}

A coarse-grained model of transcription and translation is used to describe gene
regulatory networks. The shifted Hill function captures a non-zero basal
expression rate together with a regulated expression rate that is higher
(activation) or lower (repression) than the basal level. The derivation rests on
the following assumptions.

\subsection{Assumptions}

\begin{enumerate}
    \item \emph{Quasi-steady-state binding:} DNA--TF binding and unbinding occur
    much faster than transcription and translation.
    \item \emph{TF in excess:} the transcription factor is abundant relative to
    the DNA, so its free concentration is approximately its total concentration.
\end{enumerate}

\subsection{Single-regulator model}

For a single regulator (TF), the reactions producing protein $P$ through
transcription and translation of the target gene are

\begin{align}
    \ce{DNA + TF &<=>[$k_f$][$k_r$] DNA{-}TF} \label{rxn:complex} \\
    \ce{DNA &->[$k_{ub}$] P + DNA} \label{rxn:unbound} \\
    \ce{DNA{-}TF &->[$k_{b}$] P + DNA + TF} \label{rxn:bound}
\end{align}

Reaction~\eqref{rxn:complex} describes formation of the DNA--TF complex.
Reactions~\eqref{rxn:unbound} and~\eqref{rxn:bound} together account for protein
production: for an activating TF, $k_b > k_{ub}$, and for an inhibiting TF,
$k_{ub} > k_b$. The corresponding production rate is

\begin{equation}
    \frac{d[P]}{dt} = k_{ub}\,[D] + k_{b}\,[DT],
    \label{eq:rate}
\end{equation}

where $[D]$ is the concentration of free DNA and $[DT]$ that of the DNA--TF
complex. The first term on the right-hand side is the basal expression and the
second is the regulated expression (with the roles of ``basal'' and
``regulated'' set by whether the link activates or inhibits).

Applying the quasi-steady-state assumption to TF binding,

\begin{equation}
    \frac{[DT]}{[D]\,[T]} = \frac{k_f}{k_r} \equiv K_{eq},
    \label{eq:qssa}
\end{equation}

where $[T]$ is the free TF concentration. Because the TF is in excess,
$[T]$ equals the total TF concentration. Conservation of DNA gives

\begin{equation}
    [D_0] = [D] + [DT],
    \label{eq:conservation}
\end{equation}

with $[D_0]$ the (constant) total DNA concentration. Substituting
Eq.~\eqref{eq:qssa} into Eq.~\eqref{eq:conservation},

\begin{equation}
    [D_0] = [D]\,\bigl(1 + K_{eq}\,[T]\bigr)
    \quad\Longrightarrow\quad
    [D] = \frac{[D_0]}{1 + K_{eq}\,[T]}.
    \label{eq:freeDNA}
\end{equation}

Using $[DT] = [D_0] - [D]$, Eq.~\eqref{eq:rate} becomes

\begin{equation}
    \frac{d[P]}{dt}
    = k_{ub}\,[D_0]\,\frac{1}{1 + K_{eq}\,[T]}
    + k_{b}\,[D_0]\,\frac{K_{eq}\,[T]}{1 + K_{eq}\,[T]}.
    \label{eq:inter}
\end{equation}

Defining the unbound and bound occupancy fractions

\begin{equation}
    H^{-} = \frac{1}{1 + K_{eq}\,[T]},
    \qquad
    H^{+} = \frac{K_{eq}\,[T]}{1 + K_{eq}\,[T]},
    \qquad H^{-} + H^{+} = 1,
    \label{eq:Hpm}
\end{equation}

Eq.~\eqref{eq:inter} reads

\begin{equation}
    \frac{d[P]}{dt}
    = k_{ub}\,[D_0]\,H^{-} + k_{b}\,[D_0]\,H^{+}
    = k_{ub}\,[D_0]\,\bigl(H^{-} + \lambda\,H^{+}\bigr),
    \label{eq:preNorm}
\end{equation}

where $\lambda \equiv k_b / k_{ub}$.

Let $g$ denote the maximal protein production rate. For an inhibitory link the
maximum is the unbound rate, $g = k_{ub}\,[D_0]$; for an activating link it is the
bound rate, $g = k_{b}\,[D_0] = \lambda\,k_{ub}\,[D_0]$. Normalizing
Eq.~\eqref{eq:preNorm} by $g$ therefore yields the two branches of the shifted
Hill function,

\begin{equation}
    \left.\frac{d[P]}{dt}\right._{-}
        = g\,\bigl(H^{-} + \lambda\,H^{+}\bigr),
    \qquad (\lambda \le 1,\ \text{repression})
    \label{eq:branchRep}
\end{equation}
\begin{equation}
    \left.\frac{d[P]}{dt}\right._{+}
        = \frac{g}{\lambda}\,\bigl(H^{-} + \lambda\,H^{+}\bigr)
        = g\,\Bigl(\frac{H^{-}}{\lambda} + H^{+}\Bigr).
    \qquad (\lambda \ge 1,\ \text{activation})
    \label{eq:branchAct}
\end{equation}

The two regulatory constants entering the production term are thus

\begin{equation}
    \lambda = \frac{k_b}{k_{ub}},
    \qquad
    K_{eq} = \frac{k_f}{k_r}.
\end{equation}

Here $\lambda$ is the ratio of the transcription rate with the TF bound to that
with it unbound---possibly set by TF-mediated condensate formation or RNA
polymerase recruitment---while $K_{eq}$ is the ratio of the TF's DNA binding and
unbinding rates. The half-maximal regulator concentration (threshold) follows
from $H^{-} = H^{+} = \tfrac{1}{2}$, i.e.\ $[T]_{1/2} = 1/K_{eq}$.

\paragraph{Remark (continuity and cooperativity).}
Both branches coincide at $\lambda = 1$, where $H^{-} + H^{+} = 1$ and production
equals the maximal rate $g$ independent of $[T]$: the no-regulation limit.
Cooperative binding of $n$ TF molecules replaces $K_{eq}[T]$ by
$\bigl(K_{eq}[T]\bigr)^{n}$ in Eq.~\eqref{eq:Hpm}, giving the Hill coefficient $n$
and recovering the usual sharp (switch-like) form used in the simulations.
\section{A Boolean limit explains the noise-driven enhancement of the double-positive state}
\label{sec:boolean}
 
To isolate the mechanism by which structural ($\lambda$) noise redistributes
occupancy toward high-expression states, we analyse the toggle switch in the
Boolean limit, where the exact stationary distribution can be computed in closed
form for both the deterministic (fixed-$\lambda$) and stochastic (fluctuating-$\lambda$)
cases.
 
\subsection{The two $\lambda$ extremes are not symmetric}
\label{sec:asymmetry}
 
For an inhibitory edge, the shifted Hill function entering the production term is
\begin{equation}
  H^{S}(x,\lambda) \;=\; H^{-}(x) \;+\; \lambda\,H^{+}(x),
  \qquad H^{+}(x) \equiv 1 - H^{-}(x),
  \label{eq:hshift}
\end{equation}
which is \emph{exactly affine in $\lambda$} for fixed regulator level $x$. At the
noise amplitude $\sigma=1$ the clamped $\lambda$ walk is confined to the two
boundaries of its range, $\lambda\in\{\lambda_{\min},\lambda_{\max}\}$ with
$\lambda_{\min}\approx 0.01$ and $\lambda_{\max}=1$. Evaluating
Eq.~\eqref{eq:hshift} at these two limits exposes a structural asymmetry:
\begin{align}
  \lambda = 1: \quad & H^{S}(x,1) = H^{-}(x)+H^{+}(x) = 1
    && \text{(independent of } x\text{)},\\
  \lambda \approx 0.01: \quad & H^{S}(x,0.01) \approx H^{-}(x)
    && \text{(depends on } x\text{)}.
\end{align}
The $\lambda=1$ boundary abolishes regulation \emph{unconditionally}: the target
is produced at its full rate regardless of the regulator's expression. The
$\lambda\approx 0$ boundary imposes repression only \emph{conditionally}: it
lowers the target's production rate only when the regulator $x$ is high, and has
essentially no effect when the regulator is low ($H^{-}(x)\!\approx\!1$).
An excursion toward de-repression therefore always ``fires,'' whereas an
excursion toward repression fires only when the partner node happens to be
expressed. This unconditional-up / conditional-down asymmetry is the microscopic
origin of the enhancement derived below.
 
\subsection{Boolean reduction}
 
We coarse-grain expression to $(A,B)\in\{0,1\}^2$ and each edge's regulatory
strength to $\lambda_{1},\lambda_{2}\in\{0,1\}$, where $\lambda=1$ denotes the
unconditional (de-repressed) branch and $\lambda=0$ the conditional (repressive)
branch. In the Boolean limit the fixed-$\lambda$ update rule is
\begin{equation}
  A' = \begin{cases}\overline{B} & \lambda_{1}=0\\[2pt] 1 & \lambda_{1}=1\end{cases},
  \qquad
  B' = \begin{cases}\overline{A} & \lambda_{2}=0\\[2pt] 1 & \lambda_{2}=1\end{cases}.
  \label{eq:boolupdate}
\end{equation}
Each of the four $\lambda$-configurations defines a deterministic map on the four
states $(00,01,10,11)$, with transition matrices
\begin{equation}
  M_{(0,0)}=\!\begin{pmatrix}0&0&0&1\\0&1&0&0\\0&0&1&0\\1&0&0&0\end{pmatrix}\!,\;
  M_{(1,0)}=\!\begin{pmatrix}0&0&0&1\\0&0&0&1\\0&0&1&0\\0&0&1&0\end{pmatrix}\!,\;
  M_{(0,1)}=\!\begin{pmatrix}0&0&0&1\\0&1&0&0\\0&0&0&1\\0&1&0&0\end{pmatrix}\!,\;
  M_{(1,1)}=\!\begin{pmatrix}0&0&0&1\\0&0&0&1\\0&0&0&1\\0&0&0&1\end{pmatrix}\!,
  \label{eq:configmaps}
\end{equation}
in row/column order $(00,01,10,11)$.
 
\subsection{Deterministic case ($\sigma=1$, $\lambda$ drawn once and held fixed)}
\label{sec:quenched}
 
In the deterministic (quenched) picture each parameter set draws a single
$(\lambda_{1},\lambda_{2})\in\{0.01,1\}^2$ and integrates to its attractor(s).
The four configurations yield:
\begin{center}
\begin{tabular}{clc}
\toprule
$(\lambda_1,\lambda_2)$ & Effective dynamics & Attractor(s) \\
\midrule
$(1,1)$ & $A,B$ both unconditionally on & $\{11\}$ \\
$(1,0)$ & $A$ on; $B=\overline{A}$ & $\{10\}$ \\
$(0,1)$ & $B$ on; $A=\overline{B}$ & $\{01\}$ \\
$(0,0)$ & mutual inhibition (toggle) & $\{10,\,01\}$ \\
\bottomrule
\end{tabular}
\end{center}
The $(0,0)$ configuration is the bistable toggle, whose stable attractors are the
two single-positive fixed points $\{10,01\}$; the states $00$ and $11$ are
unstable and carry no basin. Weighting the four configurations equally (each
boundary combination is equiprobable at $\sigma=1$) and splitting the bistable
config's weight symmetrically between its two attractors gives the deterministic
stationary distribution
\begin{equation}
  \pi^{\det}_{00}=0,\qquad
  \pi^{\det}_{01}=\tfrac{3}{8},\qquad
  \pi^{\det}_{10}=\tfrac{3}{8},\qquad
  \boxed{\;\pi^{\det}_{11}=\tfrac{1}{4}=0.25\;}
  \label{eq:pidet}
\end{equation}
The double-positive state $11$ is reached only under the single $(1,1)$
configuration; the conclusion $\pi^{\det}_{11}=1/4$ is independent of how the
bistable config's basins are partitioned, since $(0,0)$ never resolves to $11$.
 
\subsection{Stochastic case ($\lambda$ fluctuating)}
\label{sec:annealed}
 
When $\lambda$ fluctuates rapidly relative to the state dynamics, the effective
one-step operator is the \emph{$\lambda$-average} of the four deterministic maps
in Eq.~\eqref{eq:configmaps}:
\begin{equation}
  T \;=\; \tfrac{1}{4}\bigl(M_{(0,0)}+M_{(1,0)}+M_{(0,1)}+M_{(1,1)}\bigr)
    \;=\;
    \begin{pmatrix}
      0 & 0 & 0 & 1\\[2pt]
      0 & \tfrac{1}{2} & 0 & \tfrac{1}{2}\\[2pt]
      0 & 0 & \tfrac{1}{2} & \tfrac{1}{2}\\[2pt]
      \tfrac{1}{4} & \tfrac{1}{4} & \tfrac{1}{4} & \tfrac{1}{4}
    \end{pmatrix}.
  \label{eq:T}
\end{equation}
The chain is irreducible and aperiodic ($T_{11,11}=\tfrac14>0$), so a unique
stationary distribution exists and equals the long-run occupancy. Solving
$\pi=\pi T$ with $\sum_i\pi_i=1$ gives $\pi_{00}=\tfrac14\pi_{11}$,
$\pi_{01}=\pi_{10}=\tfrac12\pi_{11}$, hence $\tfrac{9}{4}\pi_{11}=1$ and
\begin{equation}
  \pi_{00}=\tfrac{1}{9},\qquad
  \pi_{01}=\tfrac{2}{9},\qquad
  \pi_{10}=\tfrac{2}{9},\qquad
  \boxed{\;\pi_{11}=\tfrac{4}{9}\approx 0.444\;}
  \label{eq:piann}
\end{equation}
The mean residence time in the double-positive state is
$\mathrm{MRT}(11)=\pi_{11}=4/9$.
 
\subsection{Every state flows toward $11$ under fluctuating $\lambda$}
\label{sec:column}
 
The enhancement is read directly off the last column of $T$
(Eq.~\eqref{eq:T}): the one-step probability of entering $11$ is positive from
\emph{every} state,
\begin{equation}
  T_{\,\cdot\,\to 11} = \bigl(\,1,\ \tfrac12,\ \tfrac12,\ \tfrac14\,\bigr)^{\!\top}.
\end{equation}
In particular the fully-repressed state $00$ maps to $11$ with probability~$1$:
under the update rule~\eqref{eq:boolupdate}, when neither node is expressed both
targets are de-repressed under \emph{all four} $\lambda$-configurations (the
$\lambda=1$ branch sets the target high unconditionally, and the $\lambda=0$
branch sets it high because the repressor is absent), so $00\!\to\!11$ is
$\lambda$-independent. By contrast, in the deterministic case $11$ is reachable
only under the isolated $(1,1)$ configuration. Fluctuating $\lambda$ thus lets a
parameter set that would deterministically settle into a single-positive fixed
point ($01$ or $10$) nonetheless accumulate occupancy in $11$, because every
re-draw of $\lambda$ offers a fresh, largely unconditional route upward.
 
\subsection{Comparison}
 
\begin{center}
\begin{tabular}{lccc}
\toprule
State & Deterministic $\pi^{\det}$ & Stochastic $\pi$ & Change \\
\midrule
$00$ & $0$            & $1/9\approx0.111$ & --- \\
$01$ & $3/8 = 0.375$  & $2/9\approx0.222$ & $\downarrow$ \\
$10$ & $3/8 = 0.375$  & $2/9\approx0.222$ & $\downarrow$ \\
$11$ & $1/4 = 0.250$  & $4/9\approx0.444$ & $\uparrow\ (1.78\times)$ \\
\bottomrule
\end{tabular}
\end{center}
 
Fluctuating $\lambda$ increases the occupancy of the double-positive state by a
factor of $16/9\approx1.78$ relative to the deterministic prediction from the
same clamped $\lambda$ range, drawing the redistributed mass from the two
single-positive states. This is a purely structural effect: the two calculations
use an identical, symmetric $\lambda$ ensemble confined to the same boundaries,
and differ only in whether $\lambda$ is quenched (Eq.~\eqref{eq:pidet}) or
annealed (Eq.~\eqref{eq:piann}). The gain reflects the asymmetry of
Section~\ref{sec:asymmetry}---de-repression fires unconditionally while
repression fires only conditionally---so that mixing over $\lambda$ favours the
state in which both nodes are high.
 
\paragraph{Remark.} The small stationary occupancy of $00$
($\pi_{00}=1/9$) arises from the synchronous discrete update, in which the
$(0,0)$ configuration maps $11\!\to\!00$; in the continuous ODE this appears as a
brief transient rather than a stable state. The robust, ODE-relevant conclusion
is the redistribution of occupancy from the single-positive states to the
double-positive state, not the absolute occupancy of $00$.
 

\section{Static Perturbation analysis of large-scale binary cell-fate decision networks} \label{ssec:stat_pert}
In this section, we investigate the impact of static perturbations (to edge weights and density; see \S\ref{sssec:noise_def}) on phenotypic distributions, team strengths, and their interrelation. First, we qualitatively understand the effect of noise on the steady-state distribution. We will display results for one artificial and biological network each, and provide reference to the corresponding section in the supplementary material for other networks. Fig.~\ref{fig:S1} provides the frequency of (steady-state) phenotypic distributions of $\emttwosix$ (top) and $\tseight$ (bottom) for three different noise levels (legend; arranged in decreasing order of noise level) of static perturbations to edge weights (left) and density (right). The dots represent the mean of the frequencies of the state across the iterations they have occurred. The upper and lower error bars represent the $75$-th and $25$-th percentile of this distribution, respectively. We characterize the steady states by the difference between the team expressions (see \S\ref{text:team_exp}) of the bigger and the smaller teams, i.e., $\mathrm{T1-T2}$, as shown in the x-axis. The black dashed line belongs to the corresponding wild-type network. 

Our general qualitative observations are: One, the dominant state in the unperturbed network remains dominant in perturbed networks up to $\tsw = 0.17$ for biological networks and $\tsw = 0.61$ for artificial networks, below which the state with $\mathrm{T1-T2}$ in the wild-type network closer to $0$ becomes dominant. In the latter case, prediction of the exact dominant state across noise levels isn't possible as they occur non-uniformly. Two, the number of hybrid states (states with team expressions that are neither $(0,1)$ nor $(1,0)$) increases with an increase in noise level. Three, the states that occurred in unperturbed or relatively less perturbed networks continue to occur in networks with higher perturbations. Four, there is no pattern in the frequencies for a given hybrid state across changes in noise levels. The above observations are common regardless of the network-specific nuances and the differences between the artificial and biological networks. The one contrast between artificial and biological networks is that the dominant states are always the pure states in unperturbed artificial networks, whereas, in biological networks, there are some networks with dominant states that are not pure even for $\tsw$ above the thresholds mentioned above. This is probably because of the higher interconnectivity and smaller size of artificial networks than biological networks. The frequency of dominant states decreases with increasing noise in all cases except for biological networks subject to edge-weight perturbations. Another observation for perturbed networks with wild-type dominant states not being pure is that the frequency of pure states increases with increasing noise at the expense of the dominant state frequency. In conclusion, for networks with team strength above a threshold, the teaming behavior supports the unperturbed low-dimensional phenotypic distribution. In the absence of a 'strong' teaming behavior, the network tries to behave as one unit.

\subsection{Static Edge-weight Perturbations} \label{sssec:stat_ew}
In this section, we quantify our qualitative observations for edge-weight perturbations, as summarized in Fig.~\ref{fig:S2}. We obtain the mean of team strength distributions ($\ts$) from the iterations (see App.~\ref{ap:ts_jsd_dist}) and normalize them by the corresponding team strength of the unperturbed (wild-type) network ($\tsw$). To quantify $\ts/\tsw$ as a function of edge-weight noise levels, we fit the data to an empirical model defined as:
\begin{align}
    \ts(\ewnum) &\equiv \ts, \\
    \ts(\tsw,\ewnum) &= p\cdot\tsw\exp(q\cdot\ewnum) - r, \label{eq:ts_ewnum}
\end{align}
where $\ewnum$ is the numerical value of a given edge-weight range (see Eq.~\ref{eq:rw_ew}; \S\ref{tt:ew_range_num}). The team strength values obtained from the resulting fits are plotted (as black crosses) along with $\ts$ (blue dots) against the corresponding edge-weight ranges, across networks, as shown in the top-left panel in Fig.~\ref{fig:S2}. The edge-weight ranges (x-axis) are arranged in the ascending order of the noise level they introduce, i.e., the $\ewnum$. The black-dashed line is Eq.~\ref{eq:ts_ewnum} calculated over the edge-weight noise range. The fit parameters are shown in Tab.~\ref{tab:ew_param}. We emphasize that the fit parameters aren't independent of each other and hold no physical property. This demonstrates that the teams' resilience to static edge-weight perturbations is attributed to its underlying team-based topology of the network, quantified using team strength. A natural expectation is that the emergent phenotypes from teams must also be resilient to static edge-weight perturbations, with the latter increasing with team strength, and we test this below. 

The middle and bottom panels in Fig.~\ref{fig:S2} correspond to an artificial and a biological network, respectively. We use JSD to quantify the divergence of phenotypic distribution of the perturbed network from the corresponding wild-type network, with greater JSD implying greater divergence. Therefore, higher Jensen-Shannon Divergence ($\jsd$) marks a lower resilience of the network to the introduced noise.  (see \S\ref{{sssec:jsd_def}}). The top-right panel in Fig.~\ref{fig:S2} shows the JSD distributions across edge-weight noise levels. We see a general trend in the variance of the JSD distributions, increasing with increasing edge-weight noise level. We are interested in the role of team strength on JSD in the presence of perturbations. Therefore, we propose an empirical model to establish a relationship between $\jsd$ and $\ts$ for a given network, defined as:
\begin{equation} \label{eq:jsd_ts_net}
    \jsd = a\cdot\log\left(\frac{\tsw-\ts}{\ts}\right)+b,
\end{equation}
where $b=\jsd(\ts=\tsw/2)$. Here, $\jsd$ refers to the mean of the JSD distribution, unless mentioned otherwise. We fit Eq.~\ref{eq:jsd_ts_net} to $\jsd$ and $\ts/\tsw$ across edge-weight ranges (legend) as shown in the middle and bottom-left panel in Fig.~\ref{fig:S2}. The black dashed line interpolates Eq.~\ref{eq:jsd_ts_net} over the range of $\ts/\tsw$, i.e., $[0,1]$. Note that the orange dots at the extremes are not points obtained from our simulation but correspond to $\jsd(\ts = 0) = 1$ and $\jsd(\ts = 1) = 0$ according to Eq.~\ref{eq:jsd_ts_net}. Tab.~\ref{tab:ew_param} displays the obtained parameters $a$ and $b$ from LMFIT along with their 1-sigma variations for each network. As expected, $b$ decreases with an increase in $\tsw$. We see no trend in $a$ as the first term on the right of Eq.~\ref{eq:jsd_ts_net} is independent of $\tsw$ (see Eq.~\ref{eq:ts_ewnum}). Thus, we can conclude that, for a given network, the resilience of the phenotypic distributions to static edge-weight perturbations increases with increasing team strength.

To check if the models Eqs.~\ref{eq:ts_ewnum} and~\ref{eq:jsd_ts_net} can predict the dependence of $\jsd$ on $\ts$ across networks for a particular edge-weight range, Eq.~\ref{eq:ts_ewnum} is plugged into Eq.~\ref{eq:jsd_ts_net} to result the model:
\begin{equation} \label{eq:jsd_ts_ran}
    \jsd = a\cdot\log\left(\frac{1}{p\cdot\exp(q\cdot\ewnum) - r}-1\right)+b.
\end{equation}
The $\jsd$-values along with their error-bars are calculated for constant $\ewnum$s using parameter values from Tabs.~\ref{tab:ew_art_all} and~\ref{tab:ew_art_all}. These are displayed alongside $\jsd$ (obtained from our simulations) plotted against $\ts/\tsw$ for a particular edge-weight range in the middle and bottom-right panel in Fig.~\ref{fig:S2}. We emphasize that no further fitting is done here, and Eq.~\ref{eq:jsd_ts_ran} predicts our data within 2-sigma variations. While $\jsd$ decreases with an increase in $\tsw$ for artificial networks as expected, the trend is the opposite for biological networks. This observation remains unexplained.

\section{Methods}

\subsection{Jensen--Shannon divergence}\label{sssec:jsd_def}
Jensen--Shannon divergence (JSD) is a symmetric, bounded measure of similarity between two probability distributions \cite{menendez1997jensen}, derived from the KL divergence:
\begin{align}
    \jsd(P \| Q) &= \frac{\mathrm{KL}(P \| M) + \mathrm{KL}(Q \| M)}{2}, \quad M = \frac{P + Q}{2}, \label{eq:jsd_def} \\
    \mathrm{KL}(P \| M) &= \sum_{x \in X} P(x) \log_2 \frac{P(x)}{M(x)}.
\end{align}
JSD takes values in $[0, 1]$, with $\mathrm{JSD} = 0$ indicating identical distributions. It is used here to quantify the divergence of the perturbed network's phenotypic distribution from the wild-type, serving as a quantitative measure of network resilience: higher JSD implies lower resilience.

\section{Fitting and error propagation}\label{ssec:fit_meth}
LMFIT \cite{https://doi.org/10.5281/zenodo.598352} is used for model fitting. Error bars ($1\sigma$ intervals) are calculated using Monte Carlo error propagation, which does not require the assumption that the distributions under consideration are Gaussian. Parameter values obtained from LMFIT are treated as Gaussian, and random samples drawn from these distributions are propagated through the model to generate a distribution for each derived quantity. The mean, median, and $16$th and $84$th percentiles are extracted from the resulting distribution; the upper and lower error limits are defined as the $84$th percentile minus the median and the median minus the $16$th percentile, respectively.

\subsection{Static edge weight perturbation}

To complement the continuous noise simulations and characterise the steady-state sensitivity of network phenotypes to fixed changes in interaction strength, we also performed static perturbation experiments. For each simulation run, a weight for each edge was independently sampled from a uniform distribution within one of the following ranges: $[0, 0.25]$, $[0, 0.5]$, $[0, 0.75]$, $[0, 1]$, $[0.25, 1]$, $[0.5, 1]$, $[0.75, 1]$, and $[0.95, 1]$, ordered from highest to lowest noise level. The topology and sign of interactions were held constant across all perturbations.

\noindent\textit{Assigning numerical values to edge-weight ranges:} As the smallest interval considered is $0.05$, the range $[0, 0.05]$ is assigned value $1$. The numerical value for an arbitrary range is defined as:
\begin{align}
    \ewnum([0, \mathrm{up}_{lim}]) &= \frac{\sum_{i=1}^{j} i}{j}, \quad j = \frac{\mathrm{up}_{lim}}{0.05} \nonumber \\
    \ewnum([\mathrm{low}_{lim}, 1]) &= \frac{\ewnum([0,1]) \cdot 20 - \ewnum([0, \mathrm{low}_{lim}]) \cdot k}{20 - k}, \quad k = \frac{\mathrm{low}_{lim}}{0.05} \nonumber \\
    \ewnum([\mathrm{low}_{lim}, \mathrm{up}_{lim}]) &= \frac{\ewnum([0, \mathrm{up}_{lim}]) \cdot j - \ewnum([0, \mathrm{low}_{lim}]) \cdot k}{20 - k}. \label{eq:rw_ew}
\end{align}
Note that this metric is undefined when $\mathrm{low}_{lim} = \mathrm{up}_{lim}$.


\paragraph{Initial conditions} For networks with $N < 20$, all $2^N$ possible initial states are enumerated. For $N \geq 20$, $2^{20}$ randomly generated initial states are used. In both cases, an unperturbed (wild-type) simulation with edge-weight moduli of $1$ is run as a baseline against which phenotypic distributions under perturbation are compared.

\section{Distribution of $\lambda$ under additive noise}\label{AdditiveStationary}
\subsection*{Distribution of Additive noise $X \sim \mathrm{Uniform}[0,1]$, $Z \sim \mathcal{N}(0,\sigma^2)$}
\subsubsection{Notations}
\begin{align*}
    x => Current\ position \ of\ \lambda\\
    y => Next\ value\ of\ \lambda\\
    z = y-x \sim \mathcal{N}(0,\sigma^2)\\
    p(y\mid x) => Conditional\ probability\ of\ y\\
    p(y) => probability\ of\ y\\
    p_n(x) = \frac{1}{\sqrt{2\pi}\sigma}\exp^{\frac{-x^2}{2\sigma^2}}\\
    F_n(x_0) = p_n(x < x_0) = \int_{-\infty}^{x_0}p_n(x)dx
\end{align*}
%---------------------------------------------------------
\subsubsection*{Conditional distribution and continuous density}

For $y - x = z \sim \mathcal{N}(0,\sigma^2)$:
\begin{align*}
p(y \mid x) &= p_n(y - x)\big|_{0 < y < 1}
\end{align*}

\textbf{Point masses:}
\begin{align*}
p(y \mid x) &= p_n(k \leq -x)\big|_{y=0} \\
             &= p_n(k \geq 1 - x)\big|_{y=1}
\end{align*}

\textbf{Marginal density for $0 < y < 1$:}
% \begin{align*}
% p(y) &= p_n(y - x)\big|_{0 < y < 1} \\
%      &= F_n(-x)\big|_{y \leq 0} \\
%      &= 1 - F_n(1-x)\big|_{y \geq 1}
% \end{align*}

\begin{align*}
p(y)\big|_{0 < y < 1}
  &= \int_0^1 p(y \mid x)\,dx
   = \int_0^1 p_n(y-x)\,dx \\
  &= -\int_0^1 p_n(y-x)\,d(y-x) \\
  &= -\bigl[F_n(y-x)\bigr]_0^1 \\
\end{align*}

Hence: 
\begin{equation}
    p(y)\big|_{0 < y < 1} = F_n(y) - F_n(y-1)
\end{equation}
%---------------------------------------------------------
\subsubsection*{Point masses via integration by parts}

% \textbf{Aside:} The normal pdf and its CDF satisfy
% \[
%   \frac{1}{\sqrt{2\pi}\,\sigma}\,e^{-x^2/(2\sigma^2)}, \qquad
%   \frac{1}{\sqrt{2\pi}\,\sigma}, \qquad
%   p_n(0)\,e^{-1/(2\sigma^2)}
% \]

\textbf{Mass at $y = 0$:}
\begin{align*}
p(y=0)\big|_{y=0}
  &= \int_0^1 F_n(-x)\,dx
\end{align*}

\textbf{Mass at $y = 1$:}
\begin{align*}
p(y=1)\big|_{y=1}
  &= \int_0^1 \bigl(1 - F_n(1-x)\bigr)\,dx
\end{align*}

\textbf{Substitution} $u = 1-x$, $du = -dx$:
\begin{align*}
  &= -\int_1^0 \bigl(1 - F_n(u)\bigr)\,du
   = -\int_1^0 F_n(-u)\,du
   = \int_0^1 F_n(-u)\,du \\
  &= p(y=0)\big|_{y=1} \quad \Rightarrow \quad p(y=0) = p(y=1)
\end{align*}

\textbf{IBP:} $\int u\,dv = uv - \int v\,du$, with
$u = F_n(-x)$, $v = x$:
\[
  dv = dx, \qquad du = -p_n(-x)\,dx = -p_n(x)\,dx
\]

\textbf{Gaussian identity:}
\[
  x\,p_n(x) = -\sigma^2\,p_n'(x)
  \;\Longrightarrow\;
  x\,p_n(-x)\,dx = -\sigma^2\,d\,p_n(-x)
  \;\Longrightarrow\;
  x\,p_n(n)\,dx = -d\,p_n(\cdot)
\]

\begin{align*}
\int_0^1 F_n(-x)\,dx
  &= \bigl[x\,F_n(-x)\bigr]_0^1 + \int_0^1 x\,p_n(-x)\,dx \\
  &= F_n(-1) + \bigl[p_n(0) - p_n(1)\bigr] \cdot \sigma^2
\end{align*}

Hence:
\begin{equation}
    % \begin{align*}
    p(y=0) = p(y=1) = F_n(-1) + \sigma^2\bigl[p_n(0) - p_n(1)\bigr]
    % \end{align*}
\end{equation}


%---------------------------------------------------------
\subsubsection*{CDF for $0 \leq y_0 \leq 1$}

% \textbf{Summary so far:}
% \[
% \boxed{
% \begin{aligned}
% P(0) = P(1) &= F_n(-1) + \sigma^2\bigl[p_n(0) - p_n(1)\bigr] \\
% p(y)\big|_{0<y<1} &= F_n(y) - F_n(y-1)
% \end{aligned}
% }
% \]

Note: $F_n(x_0) = P_n(x \leq x_0)$.

\textbf{CDF:}
\begin{align*}
p(y < y_0)\\
  &= p(y = 0) + p(0 < y < y_0)\\
  &= F_n(-1) + \sigma^2\bigl[p_n(0) - p_n(1)\bigr]
   + \int_0^{y_0}\bigl(F_n(y) - F_n(y-1)\bigr)\,dy
\end{align*}

\textbf{Antiderivative of $F_n$} (IBP):
\begin{align*}
\int F_n(y)\,dy
  &= F_n(y)\cdot y - \int y\,p_n(y)\,dy \\
  &= F_n(y)\,y + \sigma^2\,p_n(y) + C
\end{align*}

\textbf{Definite integral:}
\begin{align*}
\int_0^{y_0} F_n(y)\,dy
  &= y_0\,F_n(y_0) + \sigma^2\bigl[p_n(y_0) - p_n(0)\bigr]
\end{align*}

\textbf{Shifted integral} via $u = y-1$, $du = dy$, $y=0 \Rightarrow u=-1$, $y=y_0 \Rightarrow u=y_0-1$:
\begin{align*}
\int_0^{y_0} F_n(y-1)\,dy
  &= \int_{-1}^{y_0-1} F_n(u)\,du
   = \int_{-1}^{0} F_n(u)\,du + \int_0^{y_0} F_n(u)\,du
\end{align*}

Therefore:
\begin{align}
p(y < y_0) &= p(0)
   - \int_{-1}^{y_0-1} F_n(y)\,dy
   + \int_0^{y_0} F_n(y)\,dy
\end{align}

%---------------------------------------------------------
\subsubsection*{Evaluating the CDF explicitly}

\[
\int F_n(y)\,dy = y\,F_n(y) + \sigma^2\,p_n(y) + C
\]

\begin{align*}
p(y \leq y_0)
  &= p(0)
   - \int_{-1}^{y_0-1} F_n(y)\,dy
   + \int_0^{y_0} F_n(y)\,dy \\[6pt]
  &= p(0)
   - \bigl[(y_0-1)\,F_n(y_0-1) - \sigma^2\,p_n(y_0-1)\bigr] \\
  &\quad - \bigl[(-1)\,F_n(-1) + \sigma^2\,p_n(-1)\bigr] \\
  &\quad + y_0\,F_n(y_0) + \sigma^2\,p_n(y_0) - \sigma^2\,p_n(0) \\[6pt]
  &= F_n(-1) + \sigma^2\bigl[p_n(0) - p_n(1)\bigr] \\
  &\quad - (y_0-1)\,F_n(y_0-1) - \sigma^2\,p_n(y_0-1) \\
  &\quad - F_n(-1) + \sigma^2\,p_n(-1) \\
  &\quad + y_0\,F_n(y_0) + \sigma^2\,p_n(y_0) - \sigma^2\,p_n(0) \\[6pt]
  &= y_0\,F_n(y_0) - (y_0-1)\,F_n(y_0-1)
   + \sigma^2\bigl[p_n(y_0) - p_n(y_0-1)\bigr] \\[6pt]
  &= y_0\,F_n(y_0) + (1-y_0)\bigl(1 - F_n(1-y_0)\bigr)
   + \sigma^2\bigl[p_n(y_0) - p_n(1-y_0)\bigr]
\end{align*}

%---------------------------------------------------------
\subsubsection*{Final form and complement}

\[
p(y \leq y_0)
  = y_0\,F_n(y_0)
  + (1-y_0)\bigl(1 - F_n(1-y_0)\bigr)
  + \sigma^2\bigl[p_n(y_0) - p_n(1-y_0)\bigr]
\]

\textbf{Complement $P(y_0 \leq Y \leq 1)$:}
\begin{align*}
p(y_0 \leq Y \leq 1)
  &= 1 - p(Y < y_0) \\
  &= 1 - \Bigl[
      y_0\,F_n(y_0)
      + (1-y_0)\bigl(1-F_n(1-y_0)\bigr)
      + \sigma^2\bigl[p_n(y_0) - p_n(1-y_0)\bigr]
    \Bigr]
\end{align*}

\textbf{Substitution} $y_0 = 1 - y'$:
\begin{align*}
  &= 1 - \Bigl[
      (1-y')\,F_n(1-y')
      + y'\bigl(1 - F_n(y')\bigr)
      + \sigma^2\bigl[p_n(1-y') - p_n(y')\bigr]
    \Bigr] \\[4pt]
  &= 1 - (1-y')\,F_n(1-y') - y' + y'\,F_n(y') \\
  &\quad - \sigma^2\,p_n(1-y') + \sigma^2\,p_n(y') \\[4pt]
  &= (1-y')\bigl(1 - F_n(1-y')\bigr)
   + y'\,F_n(y')
   + \sigma^2\bigl[p_n(y') - p_n(1-y')\bigr]
\end{align*}

%---------------------------------------------------------
\subsubsection*{Summary and numerical example}

\[
\boxed{
\begin{aligned}
p(0) = p(1)
  &= F_n(-1) + \sigma^2\bigl[p_n(0) - p_n(1)\bigr] \\
p(y)\big|_{0<y<1}
  &= F_n(y) - F_n(y-1)
\end{aligned}
}
\]

\begin{equation}
p(y \leq y_0)
  = y_0\,F_n(y_0)
  + (1-y_0)\bigl(1 - F_n(1-y_0)\bigr)
  + \sigma^2\bigl[p_n(y_0) - p_n(1-y_0)\bigr]
  \label{additiveCDF}
\end{equation}

\textbf{Note:} this also equals $p(y_0 \leq Y \leq 1)$ by symmetry. \\[6pt]

% We find a reasonable strong agreement between the simuated and analytically derived distributions, as shown below. 

% \begin{figure}
%     \centering
%     \includegraphics[width=\linewidth]{AditiveNoiseFigures/prob_dist_compare.jpg}
%     \caption{\textbf{Comparision of CDFs generated from simulations and the analytical calculations} Each panel shows the cumulative distribution obtained from Eq \ref{additiveCDF} (red) and numerically calculated trajectories (green). The CDFs agree well with each other, with minor differences that should not have any strong impact on sampling lambda values.} 
%     \label{fig:noiseDist}
% \end{figure}


\section{Supplementary Figures}

\begin{figure}
    \centering
    \includegraphics[width=0.8\linewidth]{Figs/Fig1S1.jpg}
    \caption{
    \textbf{(A)} Heatmap of RACIPE steady-state expression (log$_2$) for nodes A and B, TS network, across all sampled parameter/steady-state combinations (ordered by increasing expression); color scale centered at the mean of the two nodes' thresholds. \textbf{(B)} Same design as Figure 1D, but for an activation-edge-range parameter ($\lambda \in [1,100]$, $\lambda_0 = 20$).}
    \label{fig:1S1}
\end{figure}

\begin{figure}
    \centering
    \includegraphics[width=0.8\linewidth]{Figs/Fig2S1.jpg}
    \caption{MRT vs.\ noise level, colored by discrete state, mean $\pm$ SD across parameter sets, for \textbf{(A)} TS (all parameter sets), \textbf{(B)} TSSA (monostable + tristable), \textbf{(C)} TT (monostable + tristable), and \textbf{(D)} TTSA (monostable + tristable).}
    \label{fig:2S1}
\end{figure}

\begin{figure}
    \centering
    \includegraphics[width=0.8\linewidth]{Figs/Fig3S1.jpg}
    \caption{\textbf{(A)} MRT boxplots (as in Fig.\ 3B) for TSSA, TT, and TTSA, stacked vertically. \textbf{(B)} Reachability scatter (as in Fig.\ 3C) for the TS network at $\sigma = 0.001, 0.01, 1.0$ (the noise levels not shown in the main figure), stacked vertically. \textbf{(C)} Reachability scatter for the TSSA network at $\sigma = 0.001, 0.01, 0.1, 1.0$.}
    \label{fig:3S1}
\end{figure}

\begin{figure}
    \centering
    \includegraphics[width=0.8\linewidth]{Figs/Fig4S1.jpg}
    \caption{\textbf{(A)} Resampling effect (as in Fig.\ 4B), collapsed to a single per-parameter-set difference and shown across all noise levels, faceted by state class (all-high/single-high/double-high) and network (TS, TSSA, TT, TTSA). \textbf{(B)} Same, for the stochastic effect. \textbf{(C)} Paramwise scatter of resampled-$\lambda$ deterministic vs.\ stochastic MRT of the all-high state, all four networks, at each of four noise levels (0.001, 0.01, 0.1, 1.0).}
    \label{fig:3bS1}
\end{figure}

\begin{figure}
    \centering
    \includegraphics[width=0.8\linewidth]{Figs/Fig4S2.jpg}
    \caption{All-high \textbf{(A)}, single-high \textbf{(B)}, and double-high \textbf{(C)} state MRT comparisons (as in Fig.\ 4D: original/resampled/stochastic, faceted by stability class), repeated at four noise levels (columns: 0.001, 0.01, 0.1, 1.0), for TS, TSSA, TT, TTSA. 2-node networks (TS, TSSA) have no double-high state distinct from all-high, so they are absent from that row.
}
    \label{fig:3bS2}
\end{figure}

\begin{figure}
    \centering
    \includegraphics[width=0.8\linewidth]{Figs/Fig5S1.jpg}
    \caption{\textbf{(A)} MRT vs.\ noise level (as in Fig.\ 5A) for TS, TSSA, and TT; single-high and double-high state classes. For TS and TSSA, whose fully-active state is intrinsically ``all-high'' rather than a distinct ``double-high'' state, this facet is relabeled ``double-high'' for consistency with the other panels, but shows the same all-high data. \textbf{(B)} $\lambda$ distributions (as in Fig.\ 4A) under Multiplicative noise. \textbf{(C)} Network topology diagrams for TT4, TT4SA, and TS4. \textbf{(D)} Deterministic attractor-combination frequency (as in Fig.\ 1C, but from the noise-simulation pipeline's zero-noise runs rather than RACIPE's own steady-state search) for TT4, TT4SA, and TS4.
}
    \label{fig:4S1}
\end{figure}

\begin{figure}
    \centering
    \includegraphics[width=0.8\linewidth]{Figs/Fig5S2.jpg}
    \caption{\textbf{(A, B)} MRT vs.\ noise level for TT4, TT4SA, and TS4; single-high, double-high, and all-high states, faceted by stability class, under Additive (A) and Multiplicative (B) noise. \textbf{(C)} State-wise MRT vs.\ noise level for each of TS4's six double-high states individually (Multiplicative noise), showing that the same subset of states dominates at every noise level. \textbf{(D)} Stochastic effect for the single-high state, across noise levels, for TSSA, TTSA, TT4, and TS4. \textbf{(E)} Stochastic effect for the all-high and double-high states, for TS, TT, and TT4. \textbf{(D$'$, E$'$)} Same networks/state classes as D and E, respectively, but showing the resampling effect instead of the stochastic effect.
}
    \label{fig:4S2}
\end{figure}

\begin{figure}
    \centering
    \includegraphics[width=0.8\linewidth]{Figs/Figure4S2.jpg}
    \caption{\textbf{Adjacency and influence matrices for the EMT and Gonadal networks.}
    A) EMT network (23 nodes). (i) Direct adjacency matrix showing the signed regulatory
    interactions between node pairs: red indicates activation ($+1$), blue indicates
    inhibition ($-1$), and white indicates no edge. Nodes are ordered by team membership,
    with the epithelial team in the lower-left block and the mesenchymal team in the
    upper-right block. (ii) Influence matrix computed as described in the Methods, showing
    the net cumulative effect between each pair of nodes across all path lengths up to
    $l_{\mathrm{max}} = 10$. Positive (red) values indicate net co-regulation; negative
    (blue) values indicate net counter-regulation. The block structure of the influence
    matrix confirms the two-team organisation of the network.
    B) Gonadal network (19 nodes), displayed in the same format as A, with the
    Sertoli-fate team and the Granulosa-fate team forming the two blocks.}
    \label{fig:5S1}
\end{figure}


% \begin{figure}
%     \centering
%     \includegraphics[width=0.8\linewidth]{Figs/Figure5S1.jpg}
%     \caption{\textbf{Team score distributions for artificial networks under all four
%     noise conditions.}
%     Two-dimensional density maps of team scores (E score on the $x$-axis, M score on the
%     $y$-axis) averaged across 50 artificial networks per condition. Columns show increasing
%     network density (Density: 1, 3, 5 representing low, intermediate, and high density
%     respectively) and rows show increasing network impurity (Impurity: 0, 1, 3). Each
%     panel corresponds to a fixed noise type and noise level ($\sigma = 0.05$):
%     (i) additive noise at $\sigma = 0$ (deterministic baseline),
%     (ii) additive noise at $\sigma = 0.05$,
%     (iii) multiplicative noise at $\sigma = 0.05$,
%     (iv) stationary noise at $\sigma = 0.05$.
%     Colour scale indicates mean frequency across networks, from low (dark) to high
%     (bright). Concentration of mass in the corners indicates terminal phenotype occupancy;
%     concentration near $(0.5, 0.5)$ indicates hybrid state dominance.}
%     \label{fig:5S2}
% \end{figure}

% \begin{figure}
%     \centering
%     \includegraphics[width=0.8\linewidth]{Figs/Figure5S2.jpg}
%     \caption{\textbf{Phenotypic score distributions for artificial networks across
%     density, impurity, and noise type.}
%     Violin plots of the phenotypic score distribution (defined as the difference between
%     team scores, ranging from $-1$ to $+1$) across 50 artificial networks per condition,
%     at multiple noise levels (colours as shown in the legend). Extreme values near $\pm 1$
%     indicate terminal phenotypes; values near $0$ indicate hybrid states.
%     A) Multiplicative noise (i) and stationary noise (ii) at zero impurity, varying
%     network density along the $x$-axis.
%     B) Stationary noise at density $0.3$, varying network impurity along the $x$-axis.
%     C--D) All three noise types (columns: additive, multiplicative, stationary) at
%     network density $0.3$ (C) and $0.5$ (D), varying network impurity along the $x$-axis.
%     Each violin shows the distribution of phenotypic scores pooled across all trajectories
%     and networks within a condition; width is proportional to frequency.}
%     \label{fig:5S2}
% \end{figure}

\begin{figure}
    \centering
    \includegraphics[width=0.8\linewidth]{Figs/Figure4S1.jpg}
    \caption{\textbf{Effect of multiplicative and stationary noise on phenotypic
    distributions in biological networks.}
    A) Team score density maps for the Gonadal network under additive noise at
    $\sigma = 0.05$, displayed as in Figure~\ref{fig:4}B. Left panel (``Rand'') shows
    the team score distribution sampled from stochastic simulation trajectories; right
    panel (``WT'') shows the distribution obtained from the steady states of the
    unperturbed network. B--C) Relative frequency of phenotypic scores across 1000
    simulation trajectories for the EMT network (i) and Gonadal network (ii) under
    multiplicative noise (B) and stationary noise (C), shown at the same noise levels as
    in Figure~\ref{fig:5}C. Colours indicate noise level as shown in the legend.
    Trajectories are initiated from steady states of the unperturbed network as in
    Figure~\ref{fig:5}.}
    \label{fig:5S2}
\end{figure}

\begin{figure*} 
\setlength{\tabcolsep}{2pt}
\resizebox{1.05\textwidth}{!}{
    \begin{tabular}{c}
        \includegraphics[height=0.3\textwidth,keepaspectratio]{Abhay_Lakshmi_old/Results/fig_1/fig1_ToggleSwitch_0.8.pdf} 
        \\[2pt]
        
        \includegraphics[height=0.3\textwidth]{Abhay_Lakshmi_old/Results/fig_1/fig1_EMT_RACIPE.pdf}\\
    \end{tabular} 
}
\caption{
\textbf{Steady-state frequencies of perturbed networks}: Steady-states are represented in terms of the difference in team expressions ($\mathrm{T1-T2}$) (x-axis);  The frequencies are for three levels of static perturbations (legend) to edge-weights (left) and density (right); The average frequencies (dots) and, the $25$-th (lower error bar) and the $75$-th (upper error bar) percentiles are displayed; The data in black belongs to the steady-state frequencies of the corresponding wild-type (unperturbed network); The top and bottom panels are of an artificial and a biological network, respectively; The titles in the panels display the network, the number of nodes and edges in the format ($x$N,$y$E) respectively, along with corresponding wild-type team strength (see Eq.~\ref{eq:ts_def}).
Generally, static perturbations negatively affect the dominant states of a network.
}
\label{fig:S1}
\end{figure*} 

\begin{figure*}
\setlength{\tabcolsep}{2pt}
\resizebox{1.05\textwidth}{!}{
\begin{tabular}{c c c c}

\textbf{(a)} &
{\hspace{-0.05cm}\includegraphics[width=2.83in,height=0.4\textwidth]{Abhay_Lakshmi_old/Results/ts_ew_ts.png}} &
\includegraphics[width=2.85in,height=0.4\textwidth]{Abhay_Lakshmi_old/Results/JSD_compiled.pdf} &
\textbf{(b)} \\[2pt]

{\raisebox{0.8cm}{\hspace{-0.15cm}\includegraphics[height=0.3\textwidth,keepaspectratio]{Abhay_Lakshmi_old/Results/jsd_ts_net_ts.png}}} &
{\hspace{-0.15cm}\includegraphics[width=2.9in,height=0.4\textwidth]{Abhay_Lakshmi_old/Results/jsd_ts_net_ts_fig.png}} &
\includegraphics[height=0.4\textwidth]{Abhay_Lakshmi_old/Results/jsd_ts_ew_ts_fig.png} &
{\raisebox{0.8cm}{\hspace{-0.15cm}\includegraphics[height=0.3\textwidth]{Abhay_Lakshmi_old/Results/jsd_ts_ew_ts.png}}} \\[1.8pt]

{\raisebox{0.8cm}{\hspace{-0.15cm}\includegraphics[height=0.3\textwidth,keepaspectratio]{Abhay_Lakshmi_old/Results/jsd_ts_net_bio.png}}} &
{\hspace{-0.05cm}\includegraphics[height=0.4\textwidth]{Abhay_Lakshmi_old/Results/jsd_ts_net_bio_fig.png}} &
\includegraphics[height=0.4\textwidth]{Abhay_Lakshmi_old/Results/jsd_ts_ew_bio_fig.png} &
{\raisebox{0.8cm}{\hspace{-0.15cm}\includegraphics[height=0.3\textwidth]{Abhay_Lakshmi_old/Results/jsd_ts_ew_bio.png}}}

\end{tabular}
}
\caption{
\textbf{Quantification $\jsd$ as a function of $\ts/\tsw$ for static edge-weight perturbations}: Unless mentioned otherwise, the network along with its nodes, edges, and wild-type team strength is shown in the title. \textit{Top:} (right) Quantification of the mean of team-strength distributions ($\ts$) normalized by the corresponding wild-type team strength ($\tsw$), $\ts/\tsw$, as a function of the edge-weight noise levels ($\ewnum$; numerical values of the edge-weight noise ranges; see \S\ref{tt:ew_range_num}) using Eq.~\ref{eq:ts_ewnum} (black-dashed line) displayed along with the error bars (see Tab.~\ref{tab:ew_param}) for $\tsten$. (left) JSD distributions for edge-weight noise ranges along with their mean ($\jsd$) (green square) and median (orange line) for $\tseight$. \textit{Middle:} (right) Quantification $\jsd$ as a function of $\ts/\tsw$ for a given artificial network: Plot of $\jsd$ and $\ts/\tsw$ across edge-weight noise ranges (legend; left) along with their fit of the Eq.~\ref{eq:jsd_ts_net} (black-dashed line; see Tab.~\ref{tab:ew_param}) for $\tseight$; For an artificial network, the resilience of steady-state distribution increases with an increase in the strength of teaming interactions. (left) $\jsd$ vs. $\ts/\tsw$ for a given edge-weight noise level (in terms of range; here $[0.75,1]$; see title) across networks (represented in terms of their wild-type team strength $\tsw$; see legend); The data in black are $\jsd$ and their error bars obtained from Eq.~\ref{eq:jsd_ts_ran} with no fitting; For a static edge-weight noise level, the resilience increases with an increase in team strength across artificial networks. \textit{Bottom:} The plots represent the same as the middle panels except for biological networks. To rule out over-fitting, we have tested our models for the edge-weight noise range $[0.2,0.3]$ (see legend of left panel). The contrast is in the right panel where $\jsd$ increases with increasing $\ts/\tsw$; This observation remains to be explained.
}
\label{fig:S2}
\end{figure*} 

\newpage
\section{Supplementary Tables}

\begin{table}[h!]
\centering
\caption{\normalsize{Threshold values used to discretize node expressions for different states in small cell-fate network motifs}}
\label{tb:Thresholds}
\begin{tabular}{|c|c|c|}
    \hline
    \textbf{Node} &  \textbf{Threshold} & \textbf{Network}\\
    \hline\hline
    A & 15.67 & Toggle Switch\\
    \hline
    B & 17.12 & Toggle Switch\\
    \hline
    
\end{tabular}


\end{table}


\noindent
\refstepcounter{table}
\label{tb:Attractors}
\begin{center}
  \small{\textbf{Table~\thetable: }}Attractor frequencies across all network motifs.
  Each sub-panel lists attractors (monostable or multistable, written as
  state-tuples joined by ``\_'') and their frequencies in RACIPE.
\end{center}

\medskip

%% ── Row 1: six 2-node networks ──────────────────────────────────────────
\noindent
\begin{minipage}[t]{0.45\linewidth}
\begin{tabular}{@{}lr@{}}
\multicolumn{2}{@{}l}{\textbf{(a) Toggle Switch (TS)}}\\
\toprule
Attractor & Freq.\\
\midrule
(0,1)          & 0.3744 \\
(1,0)          & 0.3318 \\
(0,1)\_(1,0)   & 0.1858 \\
(1,1)          & 0.0573 \\
\bottomrule
\end{tabular}
\end{minipage}
\hfill
\begin{minipage}[t]{0.45\linewidth}
\begin{tabular}{@{}lr@{}}
\multicolumn{2}{@{}l}{\textbf{(b) TS Self-Activation (TSSA)}}\\
\toprule
Attractor & Freq.\\
\midrule
(0,1)\_(1,0)           & 0.4865 \\
(0,1)                  & 0.1357 \\
(1,0)                  & 0.1243 \\
(0,0)\_(0,1)\_(1,0)    & 0.0682 \\
(0,1)\_(1,0)\_(1,1)    & 0.0369 \\
(1,0)\_(1,1)           & 0.0284 \\
(0,1)\_(1,1)           & 0.0254 \\
(0,0)\_(0,1)           & 0.0187 \\
(0,0)\_(1,0)           & 0.0184 \\
\bottomrule
\end{tabular}
\end{minipage}
\medskip

%% ── Row 2: Toggle Triad and NF Self-Activation ──────────────────────────
\noindent
\begin{minipage}[t]{0.47\linewidth}
\begin{tabular}{@{}lr@{}}
\multicolumn{2}{@{}l}{\textbf{(g) Toggle Triad (TT)}}\\
\toprule
Attractor & Freq.\\
\midrule
(0,0,1)                                    & 0.1312 \\
(0,1,0)                                    & 0.1241 \\
(1,0,0)                                    & 0.1184 \\
(0,0,1)\_(0,1,0)                           & 0.0694 \\
(0,1,0)\_(1,0,0)                           & 0.0638 \\
(0,0,1)\_(1,0,0)                           & 0.0632 \\
(1,0,1)                                    & 0.0567 \\
(1,1,0)                                    & 0.0559 \\
(0,1,1)                                    & 0.0535 \\
(0,0,1)\_(0,1,0)\_(1,0,0)                 & 0.0206 \\
(0,1,0)\_(1,0,1)                           & 0.0187 \\
(1,0,0)\_(1,0,1)                           & 0.0181 \\
(0,1,0)\_(0,1,1)                           & 0.0171 \\
(1,0,0)\_(1,1,0)                           & 0.0169 \\
(0,1,1)\_(1,0,0)                           & 0.0166 \\
(0,0,1)\_(1,1,0)                           & 0.0161 \\
(0,1,0)\_(1,1,0)                           & 0.0146 \\
(0,0,1)\_(1,0,1)                           & 0.0138 \\
(0,0,1)\_(0,1,1)                           & 0.0136 \\
(0,1,1)\_(1,1,0)                           & 0.0078 \\
(1,0,1)\_(1,1,0)                           & 0.0078 \\
(0,1,1)\_(1,0,1)                           & 0.0076 \\
(1,0,1)\_(1,0,1)                           & 0.0055 \\
(1,1,0)\_(1,1,0)                           & 0.0042 \\
(0,1,0)\_(1,0,0)\_(1,0,1)                 & 0.0039 \\
(0,0,1)\_(1,0,0)\_(1,1,0)                 & 0.0037 \\
(0,1,1)\_(0,1,1)                           & 0.0034 \\
(0,0,1)\_(0,1,0)\_(1,1,0)                 & 0.0028 \\
\bottomrule
\end{tabular}
\end{minipage}
\hfill
\begin{minipage}[t]{0.45\linewidth}
\begin{tabular}{@{}lr@{}}
\multicolumn{2}{@{}l}{\textbf{(j) Coupled Toggle Switches (TS4)}}\\
\toprule
Attractor & Freq.\\
\midrule
(0,1,0,1)\_(1,0,1,0)                               & 0.5244 \\
(1,0,1,0)                                          & 0.1313 \\
(0,1,0,1)                                          & 0.1251 \\
(0,1,0,1)\_(1,0,1,1)                               & 0.0112 \\
(1,0,1,0)\_(1,1,0,1)                               & 0.0108 \\
(0,1,0,1)\_(1,1,1,0)                               & 0.0096 \\
(0,1,1,1)\_(1,0,1,0)                               & 0.0096 \\
(0,1,0,1)\_(1,0,1,0)\_(1,1,0,0)                   & 0.0093 \\
(0,0,1,1)\_(0,1,0,1)\_(1,0,1,0)                   & 0.0092 \\
(0,1,0,1)\_(1,0,0,1)\_(1,0,1,0)                   & 0.0092 \\
(0,1,0,1)\_(0,1,1,0)\_(1,0,1,0)                   & 0.0088 \\
(0,1,0,1)\_(0,1,1,1)\_(1,0,1,0)                   & 0.0079 \\
(0,1,0,1)\_(1,0,1,0)\_(1,0,1,1)                   & 0.0065 \\
(0,1,0,1)\_(1,0,1,0)\_(1,1,1,0)                   & 0.0063 \\
(0,1,0,1)\_(1,0,1,0)\_(1,1,0,1)                   & 0.0059 \\
(0,1,0,1)\_(1,0,0,0)                               & 0.0049 \\
(0,1,0,0)\_(1,0,1,0)                               & 0.0042 \\
(0,0,1,0)\_(0,1,0,1)                               & 0.0041 \\
(0,1,0,1)\_(0,1,1,0)                               & 0.0040 \\
(0,1,0,1)\_(1,0,0,1)                               & 0.0040 \\
(0,1,0,1)\_(1,1,0,0)                               & 0.0040 \\
(1,0,0,1)\_(1,0,1,0)                               & 0.0040 \\
(0,1,0,1)\_(1,0,0,0)\_(1,0,1,0)                   & 0.0038 \\
(0,0,0,1)\_(1,0,1,0)                               & 0.0037 \\
(0,1,1,0)\_(1,0,1,0)                               & 0.0037 \\
(1,0,1,0)\_(1,1,0,0)                               & 0.0037 \\
(0,0,1,1)\_(1,0,1,0)                               & 0.0036 \\
(0,1,0,1)\_(0,1,0,1)\_(1,0,1,0)                   & 0.0035 \\
(0,0,1,1)\_(0,1,0,1)                               & 0.0034 \\
(0,0,0,1)\_(0,1,0,1)\_(1,0,1,0)                   & 0.0033 \\
(0,1,0,0)\_(0,1,0,1)\_(1,0,1,0)                   & 0.0032 \\
(1,0,1,0)\_(1,1,1,0)                               & 0.0032 \\
\bottomrule
\end{tabular}
\end{minipage}



\bigskip

%% ── Row 4: TTSA across three columns ─────────────────────────────────────
\noindent
\begin{minipage}[t]{\linewidth}
\begin{tabular}{@{}lr@{}}
\multicolumn{2}{@{}l}{\textbf{(i) Toggle Triad SA (TTSA) — cont'd across columns}}\\
\toprule
Attractor & Freq.\\
\midrule
(0,0,1)\_(0,1,0)\_(1,0,0)                         & 0.1449 \\
(0,0,1)\_(0,1,0)                                   & 0.0459 \\
(0,0,1)\_(1,0,0)                                   & 0.0444 \\
(0,1,0)\_(1,0,0)                                   & 0.0444 \\
(1,0,0)                                            & 0.0362 \\
(0,0,1)                                            & 0.0359 \\
(0,1,0)                                            & 0.0330 \\
(0,0,1)\_(0,1,0)\_(1,0,0)\_(1,0,1)                & 0.0228 \\
(1,0,0)\_(1,0,1)                                   & 0.0227 \\
(0,0,1)\_(0,1,0)\_(0,1,1)\_(1,0,0)                & 0.0220 \\
(0,0,1)\_(0,1,1)                                   & 0.0215 \\
(0,1,0)\_(0,1,1)                                   & 0.0213 \\
(0,1,0)\_(1,1,0)                                   & 0.0203 \\
(0,0,1)\_(0,1,0)\_(1,0,0)\_(1,1,0)                & 0.0201 \\
(0,0,1)\_(1,0,1)                                   & 0.0198 \\
(1,0,0)\_(1,1,0)                                   & 0.0195 \\
(0,1,1)                                            & 0.0188 \\
(1,0,1)                                            & 0.0183 \\
(0,0,1)\_(0,1,0)\_(0,1,1)                         & 0.0166 \\
(1,1,0)                                            & 0.0161 \\
(0,0,1)\_(1,0,0)\_(1,0,1)                         & 0.0140 \\
(0,1,0)\_(1,0,0)\_(1,1,0)                         & 0.0138 \\
(1,0,0)\_(1,0,1)\_(1,1,0)                         & 0.0135 \\
(0,1,0)\_(0,1,1)\_(1,1,0)                         & 0.0132 \\
(0,0,1)\_(0,1,1)\_(1,0,1)                         & 0.0117 \\
(1,0,1)\_(1,1,0)                                   & 0.0116 \\
(0,0,1)\_(0,1,0)\_(1,0,1)                         & 0.0111 \\
(0,0,1)\_(0,1,1)\_(1,0,0)                         & 0.0110 \\
(0,1,1)\_(1,1,0)                                   & 0.0102 \\
(0,0,1)\_(0,1,0)\_(1,1,0)                         & 0.0101 \\
(0,1,1)\_(1,0,1)                                   & 0.0100 \\
(0,1,0)\_(0,1,1)\_(1,0,0)                         & 0.0098 \\
(0,1,0)\_(1,0,0)\_(1,0,1)                         & 0.0097 \\
(0,0,1)\_(1,0,0)\_(1,1,0)                         & 0.0096 \\
\bottomrule
\end{tabular}
\end{minipage}
\hfill
\begin{minipage}[t]{\linewidth}
\begin{tabular}{@{}lr@{}}
\toprule
Attractor & Freq.\\
\midrule
(0,0,1)\_(0,1,0)\_(0,1,1)\_(1,0,0)\_(1,0,1)      & 0.0073 \\
(0,1,1)\_(1,0,0)                                   & 0.0071 \\
(0,0,1)\_(0,1,0)\_(1,0,0)\_(1,0,1)\_(1,1,0)      & 0.0070 \\
(0,0,1)\_(0,1,0)\_(0,1,1)\_(1,0,0)\_(1,1,0)      & 0.0064 \\
(0,0,1)\_(1,1,0)                                   & 0.0062 \\
(0,1,1)\_(1,0,1)\_(1,1,0)                         & 0.0057 \\
(0,1,0)\_(1,0,1)                                   & 0.0055 \\
(0,0,1)\_(0,1,1)\_(1,1,0)                         & 0.0045 \\
(0,1,1)\_(1,0,0)\_(1,0,1)                         & 0.0044 \\
(0,0,1)\_(0,1,1)\_(1,0,0)\_(1,1,0)                & 0.0043 \\
(0,0,1)\_(1,0,1)\_(1,1,0)                         & 0.0043 \\
(0,1,0)\_(0,1,1)\_(1,0,0)\_(1,0,1)                & 0.0043 \\
(0,0,1)\_(0,1,0)\_(0,1,1)\_(1,0,1)                & 0.0041 \\
(0,0,1)\_(0,1,1)\_(1,0,0)\_(1,0,1)                & 0.0040 \\
(0,0,1)\_(0,1,0)\_(1,0,1)\_(1,1,0)                & 0.0035 \\
(0,1,0)\_(1,0,0)\_(1,0,1)\_(1,1,0)                & 0.0032 \\
(0,1,0)\_(1,0,1)\_(1,1,0)                         & 0.0031 \\
(0,0,1)\_(0,1,0)\_(0,1,1)\_(1,0,0)\_(1,0,1)\_(1,1,0) & 0.0030 \\
(0,1,0)\_(0,1,1)\_(1,0,1)                         & 0.0030 \\
(0,1,1)\_(1,0,0)\_(1,1,0)                         & 0.0030 \\
(0,0,1)\_(0,1,1)\_(1,0,1)\_(1,1,0)                & 0.0029 \\
(0,0,1)\_(1,0,0)\_(1,0,1)\_(1,1,0)                & 0.0029 \\
(0,1,0)\_(0,1,1)\_(1,0,1)\_(1,1,0)                & 0.0028 \\
(0,0,1)\_(0,1,0)\_(0,1,1)\_(1,1,0)                & 0.0027 \\
(0,1,0)\_(0,1,1)\_(1,0,0)\_(1,1,0)                & 0.0026 \\
(0,1,1)\_(1,0,0)\_(1,0,1)\_(1,1,0)                & 0.0026 \\
(1,1,0)\_(1,1,0)                                   & 0.0021 \\
(1,0,1)\_(1,1,1)                                   & 0.0018 \\
(0,1,0)\_(0,1,1)\_(0,1,1)                         & 0.0017 \\
(0,1,0)\_(0,1,1)\_(1,0,0)\_(1,0,1)\_(1,1,0)      & 0.0017 \\
(0,1,1)\_(0,1,1)                                   & 0.0016 \\
(1,0,1)\_(1,0,1)                                   & 0.0015 \\
(0,0,1)\_(1,0,1)\_(1,0,1)                         & 0.0014 \\
(0,1,1)\_(1,1,1)                                   & 0.0013 \\
\bottomrule
\end{tabular}
\end{minipage}
\hfill
\begin{minipage}[t]{0.45\linewidth}
\begin{tabular}{@{}lr@{}}
\toprule
Attractor & Freq.\\
\midrule
(1,0,0)\_(1,0,1)\_(1,0,1)                         & 0.0013 \\
(0,0,1)\_(0,1,1)\_(0,1,1)                         & 0.0012 \\
(0,0,1)\_(0,1,1)\_(1,0,0)\_(1,0,1)\_(1,1,0)      & 0.0012 \\
(0,1,0)\_(0,1,1)\_(1,1,0)\_(1,1,0)                & 0.0012 \\
(0,1,0)\_(1,1,0)\_(1,1,0)                         & 0.0012 \\
(0,0,1)\_(0,1,0)\_(0,1,0)\_(1,0,0)                & 0.0011 \\
(0,0,1)\_(0,1,0)\_(1,0,0)\_(1,0,0)                & 0.0011 \\
(0,0,1)\_(0,1,0)\_(1,1,0)\_(1,1,0)                & 0.0011 \\
(1,0,0)\_(1,1,0)\_(1,1,0)                         & 0.0011 \\
(1,1,0)\_(1,1,1)                                   & 0.0011 \\
(0,0,0)\_(0,0,1)\_(0,1,0)\_(1,0,0)                & 0.0010 \\
(0,1,1)\_(0,1,1)\_(1,0,1)                         & 0.0010 \\
(0,1,1)\_(0,1,1)\_(1,1,0)                         & 0.0010 \\
(0,1,1)\_(1,0,1)\_(1,1,1)                         & 0.0010 \\
(1,0,0)\_(1,0,1)\_(1,1,0)\_(1,1,1)                & 0.0010 \\
(1,0,1)\_(1,1,0)\_(1,1,0)                         & 0.0010 \\
(0,0,1)\_(0,0,1)                                   & 0.0009 \\
(0,0,1)\_(0,1,0)\_(1,0,1)\_(1,0,1)                & 0.0009 \\
(0,1,1)\_(0,1,1)\_(1,0,0)                         & 0.0009 \\
(0,1,1)\_(1,1,0)\_(1,1,1)                         & 0.0009 \\
(1,0,1)\_(1,0,1)\_(1,1,0)                         & 0.0009 \\
\bottomrule
\end{tabular}
\end{minipage}

%%%%%%%%%%%%%%%%%%%%%%%%%%%%%%%%%%%%%%%%%%%%%%%%%%%%%%%%%%%%
%%% APPENDICES
%%%%%%%%%%%%%%%%%%%%%%%%%%%%%%%%%%%%%%%%%%%%%%%%%%%%%%%%%%%%
\newpage
\begingroup
\setlength{\tabcolsep}{0.7pt} % 
\renewcommand{\arraystretch}{1.5}
\begin{table}
	\caption{Fitted parameters for Equation \ref{eq:jsd_ts_ran}}
	\label{tab:ew_param}
\begin{threeparttable}
\begin{tabular}{l|c|lll|ll}
\hline
\hline
\multicolumn{1}{c|}{Network} & \multicolumn{1}{c|}{$\ts_\mathrm{wild}$} & \multicolumn{3}{c|}{$\ts$ parameters} & \multicolumn{2}{c}{$\jsd$ Parameters}\\
\hline
\hline
\multicolumn{1}{c|}{} & \multicolumn{1}{c|}{} & \multicolumn{1}{c|}{$p$} & \multicolumn{1}{c|}{$q$}  & \multicolumn{1}{c|}{$r$} & \multicolumn{1}{c|}{$a$} & \multicolumn{1}{c}{$b$}\\
\multicolumn{1}{c|}{(1)} & \multicolumn{1}{c|}{(2)} & \multicolumn{1}{c|}{(3)} & \multicolumn{1}{c|}{(4)} & \multicolumn{1}{c|}{(5)} & \multicolumn{1}{c|}{(6)} & \multicolumn{1}{c}{(7)}\\
\hline
$\tsten$  & $0.90$ & $(8.07\pm0.95)\times10^{-3}$ & $(2.34\pm0.06)\times10^{-1}$  & $(-8.07\pm5.23)\times10^{-3}$ & $(0.00\pm0.00)\times10^{-2}$ & $0.13\pm0.00$\\
$\tstwo$  & $0.88$ & $(7.39\pm0.87)\times10^{-3}$ & $(2.38\pm0.06)\times10^{-1}$  & $(-5.19\pm4.99)\times10^{-3}$ & $(0.01\pm0.00)\times10^{-2}$ & $0.01\pm0.00$\\
$\tsfour$ & $0.86$ & $(6.51\pm0.79)\times10^{-3}$ & $(2.45\pm0.06)\times10^{-1}$  & $(-3.59\pm4.80)\times10^{-3}$ & $(0.22\pm0.01)\times10^{-2}$ & $0.01\pm0.00$\\
$\tssix$  & $0.82$ & $(4.67\pm0.52)\times10^{-3}$ & $(2.61\pm0.06)\times10^{-1}$  & $(-3.09\pm3.81)\times10^{-3}$ & $(1.36\pm0.31)\times10^{-2}$ & $0.05\pm0.01$\\
$\tseight$& $0.61$ & $(2.84\pm0.01)\times10^{-3}$ & $(2.86\pm0.01)\times10^{-1}$  & $(-0.38\pm0.64)\times10^{-3}$ & $(0.89\pm0.26)\times10^{-2}$ & $0.25\pm0.01$\\
\hline
\hline
\end{tabular}
\end{threeparttable}
\end{table}
\endgroup
\newpage


% \section{Appendix}

% \section{Evaluating the effect of additive noise against the corresponding stationary distribution}

% Additive noise leads to an increase in node expression, increasing the residence time of hybrid states ((1,1) state in toggle switch). Among the three noises tested, additive noise also has the highest variance, with a high chance of achieving $\lambda \approx 1$ - weakening the inhibitions in the toggle switch and potentially allowing the nodes to express independently. Thus, it is important to understand if increasing hybrid state MRT is due to the distribution of $\lambda$s created by the additive noise or if the noise itself is required. We thus perfomed experiments to answer the following questions

% \begin{enumerate}

%     \item Do deterministic simulations of the toggle switch parameters with $\lambda$ sampled from the distributions generated by additive noise lead to an increase in the MRT of (1, 1) state?
%     \item Does modifying the noise to ensure lambda switching between 0.01 and 1 (instead of increments by normal random variable) increase the frequency of the (1, 1) state?
% \end{enumerate}
% The original formulation of additive noise follows the equation below:

% \begin{equation}
% \begin{split}
%     \lambda(t+1) &\overset{\text{Additive}}{=} 
%         \lambda(t) + \mathcal{N}(0,\sigma)\\
%     \lambda(t+1) &= \text{clip}\bigl(\lambda(t+1),\ 0.01,\ 1\bigr)
% \end{split}
% \label{additiveNoise}
% \end{equation}

% We devised two additional forms of noise - named the Jumping noise and the Extreme noise - defined as follows:


% % \centering
% \begin{equation}
% \begin{split}
%     \lambda(t+1) &\overset{\text{Jumping}}{=} 
%     \begin{cases}
%         \lambda(t) + \sigma & \text{if } \mathcal{N}(0,\sigma) > 0 \\
%         \lambda(t) - \sigma & \text{if } \mathcal{N}(0,\sigma) < 0
%     \end{cases} \\[6pt]
%     \lambda(t+1) &= \text{clip}\bigl(\lambda(t+1),\ 0.01,\ 1\bigr)
% \end{split}
% \label{jumpingNoise}
% \end{equation}

% \begin{equation}
%     \lambda (t+1) \overset{Extreme}{=} \begin{cases}
%         1 \ if \ N(0,\sigma) > 0\\
%         0.01 \ if \ N(0,\sigma) < 0
%     \end{cases}
%     \label{extremeNoise}
% \end{equation}



% We then derived the distribution of $\lambda$ for five values of $\sigma$ : ${0.001, 0.005, 0.01, 0.05, 0.1}$ numerically and simulated toggle switch as follows:

% \begin{itemize}
%     \item Sample a subset of models $\{m_i\}$ from RACIPE and extract the $\lambda$ values
%     \item For each such $\lambda$, run 100 trajectories of $\lambda$, using the corresponding noise equation (Eq \ref{additiveNoise}, \ref{jumpingNoise}, \ref{extremeNoise}), for 10000 time steps.
%     \item Collect the values of $\lambda$ from all the 100 trajectories to use as the distribution for the corresponding noise type. 
%     \item Generate 10 new models $\{m_{i_j}\}$, each model having the same parameters as the original RACIPE model, except for the $\lambda$ values sampled from the distribution generated in the previous step.
%     \item Generate deterministic trajectories for each model starting from 100 randomly generated initial conditions sampled uniformly from the range $[0, \frac{g_A}{k_A}]$, where $g_A$ is the production rate of node $A$ and $k_A$ is the degradation rate. Run each simulation for 10000 time steps
%     \item Discretize each state along these simulations using the threshold values obtained from RACIPE. 
%     \item Evaluate the mean fraction of occurance of each discrete state across all simulations and the 10 models.
%     \item Parallely run stochastic simulations on the original models of RACIPE $\{m_i\}$, starting from the steady states obtained from RACIPE, with $\lambda$ being updated at each time point according to Eq \ref{additiveNoise}, \ref{jumpingNoise}, \ref{extremeNoise}. Evaluate mean residence times of the discrete states across these stochastic trajectories.
% \end{itemize}

% \subsection{Comparing the MRTs of stochastic and deterministic trajectories}

% \begin{figure}
%     \centering
%     \begin{subfigure}[t]{0.48\linewidth}
%         \includegraphics[width=\linewidth]{AditiveNoiseFigures/11Additive.jpg}
%         \caption{Additive noise}
%         \label{fig:additive11}
%     \end{subfigure}
%     \hfill
%     \begin{subfigure}[t]{0.48\linewidth}
%         \includegraphics[width=\linewidth]{AditiveNoiseFigures/11Jumping.jpg}
%         \caption{Jumping noise}
%         \label{fig:jumping11}
%     \end{subfigure}

%     \begin{subfigure}[t]{0.48\linewidth}
%         \includegraphics[width=\linewidth]{AditiveNoiseFigures/11Extreme.jpg}
%         \caption{Extreme noise}
%         \label{fig:extreme11}
%     \end{subfigure}
%     \caption{Boxplots showing the comparison of MRTs of the (1,1) state from deterministic (red) and stochastic (green) simulations of Additive, Jumping and Extreme noise. The mean of the each distribution is shown using a '+' symbol, while the median is shown by a black horizontal line segment within each box. The boundaries of the boxes show the IQR.}
%     \label{fig:11MRT}
% \end{figure}

% For both additive and jumping noise, the mean MRT of the (1,1) state is nearly the same for both stochastic and deterministic simulations. Two differences can be highlighted: (1) the median MRT is either 0 or 1 in deterministic simulations irrespective of $\sigma$, whereas the corresponding median for stochastic simulations increases with $\sigma$. (2) The overall distribution for deterministic simulations is extremely bimodal, whereas that for the stochastic simulations shifts to higher values with increasing $\sigma$. These observations suggest that the distribution alone can indeed increase the incidence of (1,1) state in additive noise, but the increase is not consistent. The stochasticity ensures that the MRT of (1,1) remains high. Similar observations can be seen for the jumping noise case. Interestingly, in the case of extreme noise where the $\lambda$ can either be high (1) or low (0.01), the mean MRT in deterministic simulations is always lower than that of stochastic simulations, which could be explained by the fact that given only 4 combinations of $\lambda$ are possible, only one of the combinations can reliably allow the stability of (1,1) state - where both $\lambda$ values are 1. 

% \begin{figure}
%     \centering
%     \begin{subfigure}[t]{0.48\linewidth}
%         \includegraphics[width=\linewidth]{AditiveNoiseFigures/01Additive.jpg}
%         \caption{Additive noise}
%         \label{fig:additive01}
%     \end{subfigure}
%     \hfill
%     \begin{subfigure}[t]{0.48\linewidth}
%         \includegraphics[width=\linewidth]{AditiveNoiseFigures/01Jumping.jpg}
%         \caption{Jumping noise}
%         \label{fig:jumping01}
%     \end{subfigure}

%     \begin{subfigure}[t]{0.48\linewidth}
%         \includegraphics[width=\linewidth]{AditiveNoiseFigures/01Extreme.jpg}
%         \caption{Extreme noise}
%         \label{fig:extreme01}
%     \end{subfigure}
%     \caption{Boxplots showing the comparison of MRTs of the (0,1) state from deterministic (red) and stochastic (green) simulations of Additive, Jumping and Extreme noise. The mean of the each distribution is shown using a '+' symbol, while the median is shown by a black horizontal line segment within each box. The boundaries of the boxes show the IQR.}
%     \label{fig:01MRT}
% \end{figure}

% \section{Distribution of $\lambda$ under multiplicative noise}
\label{MultiplicativeStationary}

\subsection*{Distribution under Multiplicative Noise:
  $X \sim \mathrm{Uniform}[0,1]$, $Y = X(1 + Z)$, $Z \sim \mathcal{N}(0,\sigma^2)$}

\subsubsection*{Notation}
\begin{align*}
  x &\Rightarrow \text{current position of } \lambda \\
  y &\Rightarrow \text{next value of } \lambda \\
  z &= \frac{y}{x} - 1 \sim \mathcal{N}(0,\sigma^2) \\
  p(y \mid x) &\Rightarrow \text{conditional probability of } y \\
  p(y) &\Rightarrow \text{marginal probability of } y \\
  p_n(x) &= \frac{1}{\sqrt{2\pi}\,\sigma}
            \exp\!\left(-\frac{x^2}{2\sigma^2}\right) \\
  F_n(x_0) &= P(x \leq x_0) = \int_{-\infty}^{x_0} p_n(x)\,dx
\end{align*}

We also define the auxiliary function
\begin{equation}
  g(a) \;=\; \int_a^{\infty} \frac{p_n(t)}{1+t}\,dt,
  \label{eq:gdef}
\end{equation}
which will appear throughout. Note $g'(a) = -p_n(a)/(1+a)$.

%---------------------------------------------------------
\subsubsection*{Conditional distribution given $X = x$}

Since $Y = x(1+Z)$ and the state space is clamped to $[0,1]$:
\begin{itemize}
  \item $Y = 0$ iff $x(1+Z) \leq 0$, i.e.\ $Z \leq -1$ (for $x > 0$),
  \item $Y = 1$ iff $x(1+Z) \geq 1$, i.e.\ $Z \geq 1/x - 1$,
  \item $Y = y \in (0,1)$ otherwise, with $z = y/x - 1$.
\end{itemize}

Applying the change-of-variables $z = y/x - 1$, $dz = dy/x$:

\begin{equation}
  p(y \mid x) = \begin{cases}
    F_n(-1) & y = 0 \\[4pt]
    \dfrac{1}{x}\,p_n\!\left(\dfrac{y}{x} - 1\right) & 0 < y < 1 \\[8pt]
    1 - F_n\!\left(\dfrac{1}{x} - 1\right) & y = 1
  \end{cases}
  \label{eq:conditional_mult}
\end{equation}

%---------------------------------------------------------
\subsubsection*{Point mass at $y = 0$}

The boundary condition $Y = 0$ requires $Z \leq -1$, independently of $x$:
\begin{align*}
  p(y = 0)
    &= \int_0^1 P(Y=0 \mid X=x)\,dx
     = \int_0^1 F_n(-1)\,dx
\end{align*}
\begin{equation}
  \boxed{p(y=0) = F_n(-1)}
  \label{eq:mass0_mult}
\end{equation}

\textbf{Note:} Unlike the additive case, $p(y=0)$ depends only on $\sigma$
and not on the position $x$, because the sign change requires $Z < -1$
regardless of $x$.

%---------------------------------------------------------
\subsubsection*{Point mass at $y = 1$}

\begin{align*}
  p(y=1)
    &= \int_0^1 \left[1 - F_n\!\left(\tfrac{1}{x}-1\right)\right] dx
\end{align*}

\textbf{Substitution} $u = 1/x - 1$, so $x = 1/(1+u)$,
$dx = -du/(1+u)^2$; $x=1 \Rightarrow u=0$, $x \to 0^+ \Rightarrow u \to \infty$:

\begin{align*}
  p(y=1)
    &= \int_0^{\infty} \frac{1-F_n(u)}{(1+u)^2}\,du
\end{align*}

\textbf{IBP} with $A = 1-F_n(u)$, $dB = (1+u)^{-2}\,du$,
so $B = -(1+u)^{-1}$:

\begin{align*}
  &= \left[-\frac{1-F_n(u)}{1+u}\right]_0^{\infty}
     - \int_0^{\infty} \frac{p_n(u)}{1+u}\,du \\
  &= \left(0 - \frac{1-F_n(0)}{1}\right) - g(0) \\
  &= \frac{1}{2} - g(0)
\end{align*}

\begin{equation}
  \boxed{p(y=1) = \tfrac{1}{2} - g(0)}
  \label{eq:mass1_mult}
\end{equation}

\textbf{Note:} $p(y=0) \neq p(y=1)$ in general; the symmetry of the additive
case is broken because multiplicative noise scales with position.

%---------------------------------------------------------
\subsubsection*{Continuous density for $0 < y < 1$}

\begin{align*}
  p(y)
    &= \int_0^1 \frac{1}{x}\,p_n\!\left(\frac{y}{x}-1\right)dx
\end{align*}

\textbf{Substitution} $t = y/x - 1$, so $x = y/(1+t)$,
$dx = -y/(1+t)^2\,dt$; $x=1 \Rightarrow t=y-1$,
$x \to 0^+ \Rightarrow t \to \infty$:

\begin{align*}
  p(y)
    &= \int_{y-1}^{\infty}
       \frac{1+t}{y}\,p_n(t)\cdot\frac{y}{(1+t)^2}\,dt
     = \int_{y-1}^{\infty} \frac{p_n(t)}{1+t}\,dt
\end{align*}

\begin{equation}
  \boxed{p(y) = g(y-1), \qquad 0 < y < 1}
  \label{eq:density_mult}
\end{equation}

where $g$ is defined in Eq.~\eqref{eq:gdef}.
Note that for $y \in (0,1)$ the lower limit $y-1 \in (-1,0)$, so the
integral is always well-defined. 

%---------------------------------------------------------
\subsubsection*{CDF for $0 \leq y_0 \leq 1$}

\begin{align*}
  P(Y \leq y_0)
    &= p(y=0) + \int_0^{y_0} p(y)\,dy \\
    &= F_n(-1) + \int_0^{y_0} g(y-1)\,dy
\end{align*}

\textbf{Substitution} $s = y-1$, $ds = dy$:
\begin{align*}
  \int_0^{y_0} g(y-1)\,dy
    &= \int_{-1}^{y_0-1} g(s)\,ds
\end{align*}

\textbf{Switching order of integration} (Fubini) in
$\int_{-1}^{y_0-1} \int_s^{\infty} \frac{p_n(t)}{1+t}\,dt\,ds$:

\begin{align*}
  &= \int_{-1}^{\infty} \frac{p_n(t)}{1+t}
     \int_{-1}^{\min(t,\,y_0-1)} ds\,dt \\[4pt]
  &= \int_{y_0-1}^{\infty} \frac{p_n(t)}{1+t}\cdot y_0\,dt
   + \int_{-1}^{y_0-1} \frac{p_n(t)}{1+t}\cdot(t+1)\,dt \\[4pt]
  &= y_0\,g(y_0-1)
   + \int_{-1}^{y_0-1} p_n(t)\,dt \\[4pt]
  &= y_0\,g(y_0-1) + F_n(y_0-1) - F_n(-1)
\end{align*}

Therefore:
\begin{equation}
  \boxed{P(Y \leq y_0)
    = y_0\,g(y_0 - 1) + F_n(y_0 - 1)}
  \label{eq:CDF_mult}
\end{equation}

%---------------------------------------------------------
\subsubsection*{Sanity checks}

\textbf{At $y_0 = 0$:}
\[
  P(Y \leq 0) = 0 \cdot g(-1) + F_n(-1) = F_n(-1) = p(y=0) \;\checkmark
\]

\textbf{At $y_0 = 1$:}
\begin{align*}
  P(Y < 1) &= g(0) + F_n(0) = g(0) + \tfrac{1}{2} \\
  P(Y \leq 1) &= P(Y < 1) + p(y=1)
               = g(0) + \tfrac{1}{2} + \tfrac{1}{2} - g(0) = 1 \;\checkmark
\end{align*}

%---------------------------------------------------------
\subsubsection*{Summary}

\[
\boxed{
\begin{aligned}
  p(y=0) &= F_n(-1) \\
  p(y=1) &= \tfrac{1}{2} - g(0) \\
  p(y)\big|_{0<y<1} &= g(y-1) \\
  P(Y \leq y_0) &= y_0\,g(y_0-1) + F_n(y_0-1)
\end{aligned}
}
\]

where $g(a) = \displaystyle\int_a^{\infty} \frac{p_n(t)}{1+t}\,dt$.

\textbf{Key differences from additive noise:}
\begin{enumerate}
  \item $p(y=0) \neq p(y=1)$ in general — symmetry is broken.
  \item The continuous density $g(y-1)$ involves a non-elementary integral,
        unlike $F_n(y) - F_n(y-1)$ in the additive case.
  \item The CDF $y_0\,g(y_0-1) + F_n(y_0-1)$ must be evaluated numerically
        for general $\sigma$.
\end{enumerate}
\end{document}
